\chapter{凸分析与凸优化基础}
\label{ch:convex}

% ════════════════════════════════════════════════════════════════
%  数学基础部新增章：凸分析与凸优化。
%  内容来源：Robotics_Tutorial/01_数学/30_优化理论 的 10/20/30/40 四篇消化重述。
%  覆盖：凸集与凸函数·一阶/二阶判据·凸优化问题·Lagrange 对偶与 KKT·Slater·
%        共轭与 Fenchel·proximal 与 Moreau 包络·凸算法路线速览。
%  为 \cref{ch:nlopt}（最小二乘/信赖域）与控制部 MPC/H∞ 的凸性论证打地基。
%  体例严格照 lie_theory.tex：环境/记号/TikZ/章末速查。右记号粗体、SO(3)/Exp/Log 一致。
% ════════════════════════════════════════════════════════════════

\begin{practice}[前置自测]
开始之前，先试答下列问题。若已能流畅作答，可只速读本章的速查卡与定理表；若卡住，正是本章要为你解决的。
\begin{enumerate}
  \item 一个集合 $C\subseteq\mathbb{R}^n$ “凸” 是什么意思？“凸函数” 又是什么意思？两者如何由 \emph{上图}（epigraph）联系起来？
  \item 给定可微函数 $f$，你如何\emph{仅凭一阶/二阶信息}判断它是不是凸的？$\nabla^2 f\succeq\mathbf 0$ 意味着什么几何性质？
  \item 在 MPC 里你解一个二次规划得到了一个解 $\mathbf x^\star$。你凭什么相信它是\emph{全局}最优而不是局部最优？
  \item 函数 $\|\mathbf x\|_1$ 在 $\mathbf x=\mathbf 0$ 处不可微，它的 “导数” 是什么？最优性条件 $\nabla f=\mathbf 0$ 在这里如何推广？
  \item 约束优化里的 “Lagrange 乘子” $\boldsymbol\lambda$ 除了 “凑出来满足方程” 之外，有没有物理/经济含义？KKT 条件凭什么是最优的判据？
\end{enumerate}
\end{practice}

\section*{本章目标}
学完本章，你应当能够：
\begin{enumerate}
  \item \textbf{判定}任意集合 / 函数是否凸，并用\emph{保凸运算}（交、仿射、逐点上确界、部分最小化、透视）系统地构造凸性证明，而非每次回到定义硬算；
  \item \textbf{陈述并证明}分离超平面定理，理解它作为 KKT 条件与对偶理论\emph{几何根源}的角色；
  \item \textbf{区分}凸、严格凸、$\mu$-强凸与 $L$-光滑，理解条件数 $\kappa=L/\mu$ 对一阶算法收敛率的决定性影响；
  \item \textbf{计算}常见函数的次梯度，用 Fermat 规则 $\mathbf 0\in\partial f(\mathbf x^\star)$ 判定最优性，并据此理解 $\ell_1$ 正则化产生稀疏解的数学原因；
  \item \textbf{推导}凸优化问题的 Lagrange 对偶，证明弱对偶与 Slater 条件下的强对偶，精确陈述 KKT 四组件并理解其在凸问题中的充要性与 “影子价格” 解释；
  \item \textbf{掌握}共轭函数（Legendre--Fenchel 变换）的几何意义与 Fenchel--Moreau 定理（$f^{\ast\ast}=f$），理解强凸--光滑对偶；
  \item \textbf{运用} proximal 算子与 Moreau 包络，记住五个核心闭式（软阈值、块软阈值、投影、二次、SVT），理解 $\operatorname{prox}_f=(\mathbf I+\partial f)^{-1}$；
  \item \textbf{鸟瞰}凸优化算法路线（梯度下降、Nesterov 加速、近端梯度 / FISTA、ADMM、内点法），知道在给定问题结构下该选哪一个、收敛率如何。
\end{enumerate}

\subsection*{本章知识导航}
本章是一条 “\emph{先定义什么是结构良好的优化问题，再据此判定最优、计算最优、加速求解}” 的主线。结构如下：

\begin{center}
\small
\begin{tikzpicture}[
  node distance=4mm and 7mm, >=Stealth,
  blk/.style={draw, rounded corners, align=center, font=\small, inner sep=3pt, fill=main!6, draw=main!60!black},
  grp/.style={draw, rounded corners, align=center, font=\small, inner sep=3pt, fill=second!6, draw=second!60!black},
  thr/.style={draw, rounded corners, align=center, font=\small, inner sep=3pt, fill=third!8, draw=third!60!black},
  ln/.style={->, thick, gray!60!black}]
\node[blk] (set) {凸集\\ 分离超平面};
\node[blk, right=of set] (fun) {凸函数\\ 一阶 / 二阶判据};
\node[grp, right=of fun] (sub) {次梯度\\ Fermat 规则};
\node[grp, below=8mm of set] (prob) {凸优化问题\\ LP/QP/SOCP/SDP};
\node[grp, right=of prob] (dual) {Lagrange 对偶\\ Slater / KKT};
\node[thr, right=of dual] (conj) {共轭 / Fenchel\\ 强凸-光滑对偶};
\node[thr, below=8mm of prob] (prox) {proximal\\ Moreau 包络};
\node[thr, right=of prox] (alg) {算法路线\\ GD/FISTA/ADMM/IPM};
\draw[ln] (set)--(fun); \draw[ln] (fun)--(sub);
\draw[ln] (set.south) -- ++(0,-4mm) -| (prob.north);
\draw[ln] (sub.south) -- ++(0,-4mm) -| (dual.north);
\draw[ln] (prob)--(dual); \draw[ln] (dual)--(conj);
\draw[ln] (conj.south) -- ++(0,-4mm) -| (prox.north);
\draw[ln] (prox)--(alg);
\end{tikzpicture}
\end{center}

\noindent\textbf{推荐路径}：\cref{sec:cvx-set,sec:cvx-fun} 是地基（凸集、凸函数、判据），任何读者必读；\cref{sec:cvx-subgrad}（次梯度与 Fermat 规则）、\cref{sec:cvx-strong}（强凸与光滑）是连接 “最优性” 与 “收敛率” 的桥，\cref{ch:nlopt} 反复引用；\cref{sec:cvx-problem,sec:cvx-duality,sec:cvx-kkt}（凸优化问题、对偶、KKT）是控制部 MPC、SLAM 后端 SDP 松弛的理论核心；\cref{sec:cvx-conjugate,sec:cvx-prox}（共轭、proximal）服务于 ADMM 与鲁棒估计；\cref{sec:cvx-algo}（算法路线）是工程选型的速查。带 “选读” 标记的相对内部、回收锥等拓扑细节可在需要严格证明时回看。

\subsection*{前置知识桥接}
本章只用到一元微积分（梯度 $\nabla f$、Hessian $\nabla^2 f$、Taylor 展开）与线性代数（内积 $\langle\mathbf x,\mathbf y\rangle=\mathbf x^\top\mathbf y$、对称矩阵的特征值、半正定 $\mathbf A\succeq\mathbf 0$）。这些在线性代数 \cref{ch:linalg} 与矩阵微分附录 \cref{ch:matderiv} 中均有系统补齐。我们\emph{不}假设读者学过任何优化理论——本章从零把 “什么是好的优化问题” 讲清楚。

\subsection*{如果跳过本章会怎样}
\begin{itemize}
  \item 在 \cref{ch:nlopt}（非线性优化）里，你会写出最小二乘 $\min_{\mathbf x}\tfrac12\|\mathbf e(\mathbf x)\|^2$ 与信赖域子问题，却\emph{无法判断}局部解是否全局最优，也讲不清高斯--牛顿为何在残差小时表现像 “二次模型上的一步”；
  \item 在控制部的 MPC / CBF-QP 里，你会调用 OSQP、qpOASES 等 QP 求解器，但它们的\emph{终止判据就是 KKT 残差}、它们的内部就是 ADMM / 内点法——不懂本章你只能照抄参数而不会调试 “primal infeasible” 之类报错；
  \item 在 certifiable SLAM（SE-Sync）里，旋转同步这个\emph{非凸}问题被松弛成一个\emph{半正定规划}并给出全局最优\emph{证书}，其原理正是本章的对偶理论与半正定锥。
\end{itemize}

\begin{note}
\textbf{本章记号与约定.}\;
向量 / 矩阵用粗体 $\mathbf x,\mathbf A$；内积 $\langle\mathbf x,\mathbf y\rangle=\mathbf x^\top\mathbf y$，欧氏范数 $\|\mathbf x\|=\|\mathbf x\|_2$。
扩展实值函数 $f:\mathbb{R}^n\to\mathbb{R}\cup\{+\infty\}$，有效定义域 $\operatorname{dom}f=\{\mathbf x:f(\mathbf x)<+\infty\}$。
对称矩阵集 $\mathbb{S}^n$，半正定 / 正定记 $\mathbf A\succeq\mathbf 0$ / $\mathbf A\succ\mathbf 0$；向量不等式 $\mathbf x\preceq\mathbf y$ 表示逐分量 $x_i\le y_i$。
凸集记 $C,D$，凸锥记 $K$，对偶锥 $K^\ast$；凸包 $\operatorname{conv}$、仿射包 $\operatorname{aff}$、相对内部 $\operatorname{ri}$、内部 $\operatorname{int}$、闭包 $\operatorname{cl}$。
上图 $\operatorname{epi}f$、指示函数 $\delta_C$、支撑函数 $\sigma_C$、法锥 $N_C$、到 $C$ 的投影 $\Pi_C$。
次微分 $\partial f(\mathbf x)$、共轭 $f^\ast$、proximal 算子 $\operatorname{prox}_f$、Moreau 包络 $M_f^\lambda$。
强凸参数 $\mu$、光滑常数 $L$、条件数 $\kappa=L/\mu$。这套记号与 \cref{ch:lie,ch:rigid_body} 的 $\SO(3)/\SE(3)/\mathrm{Exp}/\mathrm{Log}$、概率章的 $\mathcal N(\boldsymbol\mu,\boldsymbol\Sigma)$ 全书一致。
\end{note}

% ════════════════════════════════════════════════════════════════
\section{为什么凸分析是优化理论的地基}
\label{sec:cvx-why}

\paragraph{动机：什么样的优化问题是 “好” 的？}
机器人工程师每天都在解优化问题——SLAM 后端的非线性最小二乘、MPC 的二次规划、强化学习的策略优化。但在投入具体算法之前，有一个根本问题必须回答：\textbf{什么样的优化问题是 “好” 的？} 所谓 “好”，指的是三件可以\emph{担保}的事：局部最优\emph{就是}全局最优（不会陷在坑里）、解的存在性与唯一性有理论保证、算法的收敛速度有定量刻画。凸分析提供的，正是这些保证的数学基础。

\paragraph{反面：不凸时你只能靠运气。}
想象你在调试一个轨迹优化算法。优化器返回了一个解，但你无法确定：这是全局最优还是仅仅局部最优？换一个初始点会不会得到更好的解？算法还要多少步才足够精确？如果问题是\emph{凸}的，这些问题都有明确答案；如果不是，你只能靠经验和运气。非凸问题可能有指数多个局部最优（组合优化）、梯度法只能找到局部最优、且没有高效的全局最优性判定方法。这就是为什么机器人学中 “把问题凸化”（凸松弛、凸近似）如此重要。

\begin{insight}[凸分析是优化的语法规则]\label{ins:cvx-grammar}
凸分析提供的不是某一类特殊问题的解法，而是\emph{整个优化理论的语法规则}。它定义了什么是 “结构良好” 的问题，就像线性代数定义了什么是 “结构良好” 的方程组。掌握它，你就能\emph{从猜测走向确定}：先识别问题的凸性结构，再据此选择算法、预测收敛、判定最优。
\end{insight}

\noindent 下表把本章的核心概念与机器人应用对应起来，作为全章的 “路标”：

\begin{table}[H]\centering\small
\caption{凸分析概念与机器人应用对照}
\label{tab:cvx-applications}
\begin{tabular}{ll}
\toprule
凸分析概念 & 机器人 / 估计应用 \\
\midrule
凸集的分离定理 & KKT 条件与强对偶的几何根源 \\
次梯度 & 非光滑优化（$\ell_1$ 正则、ADMM）的微分工具 \\
强凸性 & 梯度下降线性收敛的充分条件 \\
半正定锥 $\mathbb{S}^n_+$ & certifiable SLAM 的 SDP 松弛（SE-Sync） \\
二阶锥 $\mathrm{SOC}_n$ & 摩擦锥约束、推力约束轨迹优化（火箭着陆） \\
保凸运算 & CVXPY 的 DCP（纪律凸编程）验证规则 \\
Lagrange 对偶 & 求解器终止判据、非凸问题的全局最优证书 \\
proximal 算子 & ISTA/FISTA、ADMM、Moreau 光滑化（鲁棒估计 GNC） \\
\bottomrule
\end{tabular}
\end{table}

% ════════════════════════════════════════════════════════════════
\section{凸集}
\label{sec:cvx-set}

\paragraph{动机：可行域的形状决定问题的难度。}
在解约束优化时，可行域的形状决定了问题的难度。若可行域是凸的，那么沿任意可行方向走都不会 “跳出” 可行域——这个性质使得很多算法（如投影梯度法）可以高效工作，也使得 “局部 $\Rightarrow$ 全局” 的奇迹成为可能。

\begin{definition}[凸集与凸组合]\label{def:convex-set}
集合 $C\subseteq\mathbb{R}^n$ 是\textbf{凸集}，当且仅当对所有 $\mathbf x,\mathbf y\in C$ 与所有 $\theta\in[0,1]$，有
\begin{equation}
  \theta\mathbf x+(1-\theta)\mathbf y\in C.
  \label{eq:convex-set-def}
\end{equation}
点 $\theta_1\mathbf x_1+\dots+\theta_k\mathbf x_k$（其中 $\theta_i\ge0$、$\sum_i\theta_i=1$）称为 $\mathbf x_1,\dots,\mathbf x_k$ 的\textbf{凸组合}。集合 $S$ 的\textbf{凸包} $\operatorname{conv}(S)$ 是包含 $S$ 的最小凸集，等于 $S$ 中所有点的凸组合之全体。
\end{definition}

\noindent\textbf{几何直觉.}\;凸集是 “没有凹陷的集合”：任取两点，其连线段整段仍在集合内。它像一个充了气的气球，无法从外部按出凹坑；非凸集则像一块有凹槽的橡皮泥，两点之间的直线可能穿出集合（\cref{fig:cvx-set}）。

\begin{figure}[ht]
\centering
\begin{tikzpicture}[scale=1.0,>=Stealth]
  % convex blob (left)
  \begin{scope}
    \draw[thick, main, fill=main!8] plot[smooth cycle, tension=0.8]
      coordinates {(-2.4,0) (-1.8,0.9) (-0.9,1.0) (-0.3,0.4) (-0.6,-0.7) (-1.6,-0.9) (-2.3,-0.5)};
    \fill (-1.9,0.55) circle (1.0pt) node[above left, font=\scriptsize] {$\mathbf x$};
    \fill (-0.7,-0.45) circle (1.0pt) node[below right, font=\scriptsize] {$\mathbf y$};
    \draw[second, thick] (-1.9,0.55)--(-0.7,-0.45);
    \node[font=\small, main] at (-1.4,-1.4) {凸：线段在内};
  \end{scope}
  % non-convex blob (right)
  \begin{scope}[xshift=4.6cm]
    \draw[thick, red!60!black, fill=red!6] plot[smooth cycle, tension=0.7]
      coordinates {(-2.4,0) (-1.6,0.9) (-0.9,0.2) (-0.4,0.9) (0.2,0.1) (-0.2,-0.8) (-1.4,-0.9) (-2.2,-0.4)};
    \fill (-2.0,0.45) circle (1.0pt) node[above left, font=\scriptsize] {$\mathbf x$};
    \fill (0.0,-0.2) circle (1.0pt) node[below right, font=\scriptsize] {$\mathbf y$};
    \draw[second, thick, dashed] (-2.0,0.45)--(0.0,-0.2);
    \node[font=\small, red!60!black] at (-1.1,-1.4) {非凸：线段穿出};
  \end{scope}
\end{tikzpicture}
\caption{凸集（左）与非凸集（右）。凸性是\emph{代数}性质（线性组合封闭），与边界是否光滑无关——正方形有角，仍是凸集。}
\label{fig:cvx-set}
\end{figure}

\paragraph{两类推广：仿射集与凸锥。}
放宽或加强 \cref{eq:convex-set-def} 中对 $\theta$ 的限制，得到两个重要变体。\textbf{仿射集} 要求对\emph{所有} $\theta\in\mathbb{R}$（不限于 $[0,1]$）都封闭，它是 “直线的推广”，任何仿射集可写成 $\{\mathbf x_0+\mathbf v:\mathbf v\in V\}$（$V$ 为子空间）；仿射集一定凸，反之不然。\textbf{锥} 要求对所有 $\theta\ge0$ 有 $\theta\mathbf x\in K$；既是凸集又是锥的称\textbf{凸锥}，等价于对 $\theta_1,\theta_2\ge0$ 有 $\theta_1\mathbf x_1+\theta_2\mathbf x_2\in K$。

\begin{insight}[锥规划的层级建立在凸锥包含链上]\label{ins:cone-hierarchy}
凸锥是优化理论中最重要的结构之一。整个锥规划（LP、SOCP、SDP）的层级，建立在不同凸锥的包含关系之上：非负象限 $\mathbb{R}^n_+$（对应 LP）$\subset$ 二阶锥 $\mathrm{SOC}_n$（SOCP）$\subset$ 半正定锥 $\mathbb{S}^n_+$（SDP）。识别可行域落在哪一级锥里，就等于选定了求解器的类别（见 \cref{sec:cvx-problem}）。
\end{insight}

\subsection{常备凸集家族}
\label{subsec:cvx-set-zoo}

认识一批 “常备” 凸集，就像掌握一套零件库，遇到新问题时能快速判断 “这个可行域长什么样”。

\begin{itemize}
  \item \textbf{超平面与半空间}：$\{\mathbf x:\mathbf a^\top\mathbf x=b\}$（$\mathbf a\neq\mathbf 0$）与 $\{\mathbf x:\mathbf a^\top\mathbf x\le b\}$。半空间是最基本的凸集构件——\emph{任何闭凸集都是（可能无穷多个）半空间的交}。
  \item \textbf{多面体}：有限个半空间与超平面的交 $P=\{\mathbf x:\mathbf A\mathbf x\preceq\mathbf b,\ \mathbf C\mathbf x=\mathbf d\}$。LP 的可行域就是多面体；其\emph{顶点}（不能写成集合中其他两点凸组合的点）是线性规划最优解的取值处——单纯形法的理论基础。
  \item \textbf{欧氏球与椭球}：$B(\mathbf x_c,r)=\{\mathbf x:\|\mathbf x-\mathbf x_c\|\le r\}$ 与 $\mathcal E=\{\mathbf x:(\mathbf x-\mathbf x_c)^\top\mathbf P^{-1}(\mathbf x-\mathbf x_c)\le1\}$（$\mathbf P\succ\mathbf 0$）。椭球半轴由 $\mathbf P$ 的特征向量 / 特征值的平方根决定；在估计与 MPC 中常用来表示不确定性区域。
  \item \textbf{二阶锥（Lorentz 锥）}：$\mathrm{SOC}_n=\{(\mathbf x,t)\in\mathbb{R}^{n+1}:\|\mathbf x\|\le t\}$。出现在摩擦锥约束（接触力须在摩擦锥内）、推力约束（火箭着陆）、鲁棒约束中。
  \item \textbf{半正定锥}：$\mathbb{S}^n_+=\{\mathbf X\in\mathbb{S}^n:\mathbf X\succeq\mathbf 0\}=\{\mathbf X\in\mathbb{S}^n:\mathbf v^\top\mathbf X\mathbf v\ge0,\ \forall\mathbf v\}$。出现在 certifiable SLAM、LMI 控制综合、SOS 多项式验证中。
\end{itemize}

\begin{proposition}[二阶锥与半正定锥的凸性]\label{prop:soc-psd-convex}
$\mathrm{SOC}_n$ 与 $\mathbb{S}^n_+$ 都是凸锥。
\end{proposition}
\begin{proof}
\emph{$\mathrm{SOC}_n$}：取 $(\mathbf x_1,t_1),(\mathbf x_2,t_2)\in\mathrm{SOC}_n$、$\theta\in[0,1]$，由三角不等式与范数正齐性
\[
  \|\theta\mathbf x_1+(1-\theta)\mathbf x_2\|\le\theta\|\mathbf x_1\|+(1-\theta)\|\mathbf x_2\|\le\theta t_1+(1-\theta)t_2,
\]
故凸组合仍在锥内。\emph{$\mathbb{S}^n_+$}：设 $\mathbf X_1,\mathbf X_2\succeq\mathbf 0$，对任意 $\mathbf v$，
\[
  \mathbf v^\top\!\big(\theta\mathbf X_1+(1-\theta)\mathbf X_2\big)\mathbf v=\theta\,\mathbf v^\top\mathbf X_1\mathbf v+(1-\theta)\,\mathbf v^\top\mathbf X_2\mathbf v\ge0,
\]
故 $\theta\mathbf X_1+(1-\theta)\mathbf X_2\succeq\mathbf 0$。锥性质（对 $\theta\ge0$ 的数乘封闭）由定义直接可见。
\end{proof}

\begin{pitfall}[编程陷阱：半正定判定只检查对角元]
半正定（PSD）要求\emph{所有特征值}非负，这是全局性质，\textbf{不能}从对角元推出。例如 $\big[\begin{smallmatrix}1&2\\2&1\end{smallmatrix}\big]$ 对角元全为正，但特征值为 $3,-1$，并非 PSD。正确做法是计算最小特征值（如 \texttt{SelfAdjointEigenSolver}）并检查其 $\ge0$，而不是 \texttt{A.\allowbreak diagonal().\allowbreak minCoeff()>=0}。
\end{pitfall}

\subsection{相对内部（选读但要害）}
\label{subsec:cvx-relint}

\begin{definition}[相对内部]\label{def:relint}
凸集 $C$ 的\textbf{相对内部} $\operatorname{ri}(C)$ 是 $C$ 在其仿射包 $\operatorname{aff}(C)$ 中的内部：
\begin{equation}
  \operatorname{ri}(C)=\{\mathbf x\in C:\exists r>0,\ B(\mathbf x,r)\cap\operatorname{aff}(C)\subseteq C\}.
  \label{eq:relint-def}
\end{equation}
\end{definition}

为什么需要这个概念？考虑 $\mathbb{R}^3$ 中的一条线段 $[\mathbf a,\mathbf b]$，它的\emph{内部}（在 $\mathbb{R}^3$ 意义下）是空的，但它的相对内部（在自身仿射包——一条直线——中的内部）是开线段 $(\mathbf a,\mathbf b)$。许多定理（“凸函数在 $\operatorname{ri}(\operatorname{dom}f)$ 上连续”、分离定理的强版本、Fenchel--Moreau 定理）必须用 $\operatorname{ri}$ 而非 $\operatorname{int}$ 才能正确陈述——否则低维凸集的结论全部失效。一个保证我们后面会反复用到：\emph{非空凸集的相对内部永远非空}。

\begin{table}[H]\centering\small
\caption{内部与相对内部（在 $\mathbb{R}^3$ 中观察）}
\label{tab:relint}
\begin{tabular}{lll}
\toprule
集合 & 内部 $\operatorname{int}$ & 相对内部 $\operatorname{ri}$ \\
\midrule
球体 & 开球 & 同内部 \\
圆盘（2D 嵌入 3D） & $\emptyset$ & 开圆盘 \\
线段 & $\emptyset$ & 开线段 \\
单点 & $\emptyset$ & 该点本身 \\
\bottomrule
\end{tabular}
\end{table}

\begin{practice}
\begin{enumerate}
  \item 证明两个凸集的交是凸集；给出反例说明并一般不凸。
  \item 证明单纯形 $\Delta_n=\{\mathbf x\in\mathbb{R}^n:\mathbf x\succeq\mathbf 0,\ \mathbf 1^\top\mathbf x=1\}$ 是凸集，并指出其极值点。
  \item 验证椭球 $\mathcal E$ 是单位球在仿射变换 $\mathbf x\mapsto\mathbf P^{1/2}\mathbf x+\mathbf x_c$ 下的像，从而由保凸运算（\cref{sec:cvx-separation}）说明它凸。
\end{enumerate}
\end{practice}

% ════════════════════════════════════════════════════════════════
\section{保凸运算与分离超平面定理}
\label{sec:cvx-separation}

\paragraph{动机：系统地判定凸性。}
实际问题的可行域往往不是单一的简单凸集，而由各种运算组合而来。\textbf{保凸运算}告诉我们：哪些操作施加在凸集上结果仍凸。这是系统性判断凸性的核心工具，也是 CVXPY 的 DCP 规则的数学基础。

\begin{theorem}[凸集的保凸运算]\label{thm:set-ops}
设 $C,C_1,C_2$ 为凸集，则下列结果仍为凸集：
\begin{enumerate}
  \item \textbf{交}：任意（甚至无穷）多个凸集的交 $\bigcap_\alpha C_\alpha$；
  \item \textbf{仿射像与原像}：仿射映射 $f(\mathbf x)=\mathbf A\mathbf x+\mathbf b$ 下的 $f(C)$ 与 $f^{-1}(C)$；
  \item \textbf{Minkowski 和}：$C_1+C_2=\{\mathbf x+\mathbf y:\mathbf x\in C_1,\mathbf y\in C_2\}$；
  \item \textbf{笛卡尔积}：$C_1\times C_2$；
  \item \textbf{透视像}：透视函数 $P(\mathbf x,t)=\mathbf x/t$（$t>0$）下的 $P(C)$。
\end{enumerate}
\end{theorem}
\begin{proof}
仅证 “交”（其余类似，把定义式 \cref{eq:convex-set-def} 代入即可）。取 $\mathbf x,\mathbf y\in\bigcap_\alpha C_\alpha$，则对每个 $\alpha$ 有 $\mathbf x,\mathbf y\in C_\alpha$；由 $C_\alpha$ 凸，$\theta\mathbf x+(1-\theta)\mathbf y\in C_\alpha$ 对所有 $\alpha$ 成立，故落在交集中。
\end{proof}

\begin{pitfall}[概念误区：凸集的并也是凸集]
两个凸集的\emph{并}一般\textbf{不}凸（如两个不相交的球）。只有\emph{交}保凸；但 Minkowski \emph{和} 保凸——务必区分 “并” 与 “和”。
\end{pitfall}

\subsection{分离超平面定理}
\label{subsec:cvx-sep-thm}

这是凸分析中最重要的定理之一，是 KKT 条件与对偶理论的几何根源。

\begin{theorem}[分离超平面定理]\label{thm:separation}
设 $C,D$ 是两个不相交的凸集（$C\cap D=\emptyset$）。则存在 $\mathbf a\neq\mathbf 0$ 与 $b$，使得
\begin{equation}
  \mathbf a^\top\mathbf x\le b\le\mathbf a^\top\mathbf y,\qquad\forall\mathbf x\in C,\ \forall\mathbf y\in D.
  \label{eq:separation}
\end{equation}
\end{theorem}

\noindent\textbf{几何直觉.}\;两个不相交的凸集之间一定能塞进一张超平面，把它们分到两侧（\cref{fig:cvx-sep}）。若集合不凸，这未必成立——平面上两个 “半月形” 虽不相交却可能交错嵌套，无法用一条直线分开。凸性正排除了这种嵌套：凸集的边界向外鼓，不会深入另一集合的内部。

\begin{figure}[ht]
\centering
\begin{tikzpicture}[scale=1.0,>=Stealth]
  \draw[thick, main, fill=main!8] (-2.3,0) circle (0.95);
  \node[main, font=\small] at (-2.3,0) {$C$};
  \draw[thick, second, fill=second!8] (1.9,0.2) circle (0.95);
  \node[second!60!black, font=\small] at (1.9,0.2) {$D$};
  % separating hyperplane
  \draw[very thick, third!70!black] (-0.4,-1.6) -- (-0.1,1.7);
  \node[third!70!black, font=\scriptsize, rotate=80] at (0.05,1.2) {$\mathbf a^\top\mathbf x=b$};
  % normal arrow
  \draw[->, thick, third!70!black] (-0.25,0.05) -- (0.55,0.18) node[below right, font=\scriptsize] {$\mathbf a$};
\end{tikzpicture}
\caption{分离超平面：不相交的两个凸集 $C,D$ 被法向量为 $\mathbf a$ 的超平面分到两侧。这是 KKT 条件、Lagrange 对偶、所有一阶最优性条件的几何根源。}
\label{fig:cvx-sep}
\end{figure}

\begin{derivation}[闭凸集与外点的严格分离：完整证明]
我们给出最常用、最严格的版本：设 $C$ 是\emph{闭}凸集，$\mathbf x_0\notin C$，构造一张严格分离超平面。证明分四步，后面的投影性质（变分不等式）本身在 \cref{sec:cvx-prox} 还会复用。

\textbf{第一步（最近点存在且唯一）.}\;设 $d=\inf_{\mathbf x\in C}\|\mathbf x-\mathbf x_0\|$，取点列 $\mathbf x_k\in C$ 使 $\|\mathbf x_k-\mathbf x_0\|\to d$。由平行四边形恒等式，
\[
  \|\mathbf x_k-\mathbf x_l\|^2=2\|\mathbf x_k-\mathbf x_0\|^2+2\|\mathbf x_l-\mathbf x_0\|^2-4\Big\|\tfrac{\mathbf x_k+\mathbf x_l}{2}-\mathbf x_0\Big\|^2.
\]
由 $C$ 凸，中点 $\tfrac{\mathbf x_k+\mathbf x_l}{2}\in C$，故末项 $\ge d^2$；当 $k,l$ 充分大时前两项 $\to d^2$，于是 $\|\mathbf x_k-\mathbf x_l\|^2\to0$，$\{\mathbf x_k\}$ 是 Cauchy 列，极限 $\mathbf p$ 存在且（$C$ 闭）落在 $C$ 中。唯一性同理：若有两个最近点 $\mathbf p_1\neq\mathbf p_2$，其中点离 $\mathbf x_0$ 更近，矛盾。

\textbf{第二步（投影的变分不等式特征）.}\;$\mathbf p=\Pi_C(\mathbf x_0)$ 当且仅当
\begin{equation}
  \langle\mathbf x_0-\mathbf p,\ \mathbf x-\mathbf p\rangle\le0,\qquad\forall\mathbf x\in C.
  \label{eq:proj-variational}
\end{equation}
取 $\mathbf x\in C$、$t\in[0,1]$，由凸性 $\mathbf p+t(\mathbf x-\mathbf p)\in C$，最近性给出
\[
  \|\mathbf x_0-\mathbf p\|^2\le\|\mathbf x_0-\mathbf p-t(\mathbf x-\mathbf p)\|^2=\|\mathbf x_0-\mathbf p\|^2-2t\langle\mathbf x_0-\mathbf p,\mathbf x-\mathbf p\rangle+t^2\|\mathbf x-\mathbf p\|^2,
\]
除以 $t>0$ 再令 $t\to0^+$ 即得 \cref{eq:proj-variational}。

\textbf{第三步（构造超平面）.}\;令 $\mathbf a=\mathbf x_0-\mathbf p\neq\mathbf 0$。对任意 $\mathbf x\in C$，由 \cref{eq:proj-variational} 有 $\mathbf a^\top(\mathbf x-\mathbf p)\le0$，即 $\mathbf a^\top\mathbf x\le\mathbf a^\top\mathbf p$；而 $\mathbf a^\top\mathbf x_0=\mathbf a^\top\mathbf p+\|\mathbf a\|^2>\mathbf a^\top\mathbf p$。取 $b=\mathbf a^\top\mathbf p+\tfrac12\|\mathbf a\|^2$，则超平面 $\{\mathbf x:\mathbf a^\top\mathbf x=b\}$ 严格分离 $\mathbf x_0$ 与 $C$。

\textbf{第四步（推广到两个不相交闭凸集）.}\;若 $C,D$ 不相交、闭凸且至少一个有界，令 $E=C-D=\{\mathbf x-\mathbf y:\mathbf x\in C,\mathbf y\in D\}$。由 \cref{thm:set-ops}，$E$ 凸，且 $\mathbf 0\notin E$（否则 $C\cap D\neq\emptyset$）。对 $\mathbf 0$ 与 $E$（或其闭包）应用第一至三步，所得超平面同时分离 $C$ 与 $D$。
\end{derivation}

\begin{theorem}[支撑超平面定理]\label{thm:supporting}
设 $C$ 是凸集，$\mathbf x_0\in\partial C$（边界点）。则存在 $\mathbf a\neq\mathbf 0$ 使 $\mathbf a^\top\mathbf x\le\mathbf a^\top\mathbf x_0$ 对所有 $\mathbf x\in C$。
\end{theorem}
\begin{proof}
取 $C$ 外的点列 $\mathbf y_k\to\mathbf x_0$。对每个 $\mathbf y_k$ 用 \cref{thm:separation} 得单位法向量 $\mathbf a_k$（$\|\mathbf a_k\|=1$）使 $\mathbf a_k^\top\mathbf x\le\mathbf a_k^\top\mathbf y_k$（$\forall\mathbf x\in C$）。单位球面紧，取收敛子列 $\mathbf a_{k_j}\to\mathbf a^\ast$（$\|\mathbf a^\ast\|=1$）。对任意 $\mathbf x\in C$，由连续性 $\mathbf a^{\ast\top}\mathbf x=\lim_j\mathbf a_{k_j}^\top\mathbf x\le\lim_j\mathbf a_{k_j}^\top\mathbf y_{k_j}=\mathbf a^{\ast\top}\mathbf x_0$。
\end{proof}

\begin{insight}[分离定理的本质]\label{ins:sep-essence}
分离超平面定理的本质是：凸集在其边界上不会 “缠绕” 或 “自交”，于是我们总能用线性不等式（超平面）从外部 “逼近” 凸集。这正是线性规划对偶、Lagrange 乘子法、以及一切一阶最优性条件的几何根源——后面 \cref{sec:cvx-duality} 的强对偶证明、\cref{sec:cvx-conjugate} 的 $f^{\ast\ast}=f$ 都直接踩在它上面。
\end{insight}

\subsection{推论：支撑函数、对偶锥与闭凸集的半空间表示}
\label{subsec:cvx-sep-cor}

\begin{itemize}
  \item \textbf{闭凸集是半空间的交}：由 \cref{thm:separation}，任何闭凸集 $C=\bigcap\{\mathbf x:\mathbf a^\top\mathbf x\le b\}$（交遍所有包含 $C$ 的半空间）。这是多面体表示与线性规划对偶的基础。
  \item \textbf{支撑函数}：$\sigma_C(\mathbf y)=\sup_{\mathbf x\in C}\langle\mathbf y,\mathbf x\rangle$，给出 $C$ 在方向 $\mathbf y$ 上的 “最大延伸”。它永远凸、正齐次，且完全刻画闭凸集（$C_1=C_2\iff\sigma_{C_1}=\sigma_{C_2}$）；后面会看到它是指示函数的共轭（\cref{sec:cvx-conjugate}）。
  \item \textbf{对偶锥（极锥）}：$K^\ast=\{\mathbf y:\langle\mathbf y,\mathbf x\rangle\ge0,\ \forall\mathbf x\in K\}$。重要的\emph{自对偶}例子：$(\mathbb{R}^n_+)^\ast=\mathbb{R}^n_+$、$(\mathrm{SOC}_n)^\ast=\mathrm{SOC}_n$、$(\mathbb{S}^n_+)^\ast=\mathbb{S}^n_+$。自对偶性在锥规划对偶中扮演核心角色（\cref{subsec:cvx-conic}）。
\end{itemize}

\begin{practice}
\begin{enumerate}
  \item 利用保凸运算证明 $\{\mathbf X:\|\mathbf X\|_2\le1\}$（谱范数球）是凸集。
  \item 计算 $\ell_1$ 单位球 $\{\mathbf x:\|\mathbf x\|_1\le1\}$ 的支撑函数（答案应为 $\|\mathbf y\|_\infty$）。
  \item 证明半正定锥自对偶：$(\mathbb{S}^n_+)^\ast=\mathbb{S}^n_+$。
\end{enumerate}
\end{practice}

% ════════════════════════════════════════════════════════════════
\section{凸函数及其判据}
\label{sec:cvx-fun}

\paragraph{动机：四种等价定义各有所长。}
判断一个函数是否凸是优化中最基本的操作。“凸” 有多种等价刻画，每种适用于不同场景：定义式对\emph{任何}函数都成立；epigraph 把函数凸性翻译成集合凸性；一阶条件给出最优性所需的几何直觉；二阶条件用 Hessian 判定，实际计算最常用。

\begin{definition}[凸函数及其上图]\label{def:convex-fun}
$f:\mathbb{R}^n\to\mathbb{R}\cup\{+\infty\}$ 是\textbf{凸函数}，如果 $\operatorname{dom}f$ 凸，且对所有 $\mathbf x,\mathbf y\in\operatorname{dom}f$、$\theta\in[0,1]$，
\begin{equation}
  f(\theta\mathbf x+(1-\theta)\mathbf y)\le\theta f(\mathbf x)+(1-\theta)f(\mathbf y).
  \label{eq:convex-fun-def}
\end{equation}
其\textbf{上图} $\operatorname{epi}f=\{(\mathbf x,t):\mathbf x\in\operatorname{dom}f,\ t\ge f(\mathbf x)\}$。则 $f$ 凸 $\iff\operatorname{epi}f$ 是凸集。
\end{definition}

\noindent\textbf{几何直觉.}\;凸函数图形上任意两点的连线段位于图形上方或重合——函数 “碗状向上弯”。epigraph 定义把函数凸性问题转化为集合凸性问题：当我们想用集合论工具（如分离定理）研究凸函数时，就通过上图来 “翻译”。

\begin{theorem}[一阶与二阶判据]\label{thm:cvx-criteria}
设 $f$ 在凸的 $\operatorname{dom}f$ 上可微，则
\begin{equation}
  f\ \text{凸}\iff f(\mathbf y)\ge f(\mathbf x)+\nabla f(\mathbf x)^\top(\mathbf y-\mathbf x),\qquad\forall\mathbf x,\mathbf y.
  \label{eq:cvx-first-order}
\end{equation}
若 $f$ 二次可微，则
\begin{equation}
  f\ \text{凸}\iff\nabla^2 f(\mathbf x)\succeq\mathbf 0,\qquad\forall\mathbf x\in\operatorname{dom}f.
  \label{eq:cvx-second-order}
\end{equation}
\end{theorem}
\begin{proof}
\emph{一阶（$\Rightarrow$）}：由 \cref{eq:convex-fun-def}，对 $\theta\in(0,1]$ 有 $f(\mathbf x+\theta(\mathbf y-\mathbf x))-f(\mathbf x)\le\theta(f(\mathbf y)-f(\mathbf x))$；除以 $\theta$ 并令 $\theta\to0^+$，左边趋于方向导数 $\nabla f(\mathbf x)^\top(\mathbf y-\mathbf x)$，得 \cref{eq:cvx-first-order}。\emph{一阶（$\Leftarrow$）}：取 $\mathbf z=\theta\mathbf x+(1-\theta)\mathbf y$，对 $(\mathbf z,\mathbf x)$ 与 $(\mathbf z,\mathbf y)$ 分别用 \cref{eq:cvx-first-order}，按 $\theta$ 与 $1-\theta$ 加权相加，并注意 $\theta\mathbf x+(1-\theta)\mathbf y-\mathbf z=\mathbf 0$，即得 \cref{eq:convex-fun-def}。

\emph{二阶（$\Leftarrow$）}：由带积分余项的 Taylor 展开 $f(\mathbf y)=f(\mathbf x)+\nabla f(\mathbf x)^\top(\mathbf y-\mathbf x)+\tfrac12(\mathbf y-\mathbf x)^\top\nabla^2 f(\boldsymbol\xi)(\mathbf y-\mathbf x)$，末项因 $\nabla^2 f(\boldsymbol\xi)\succeq\mathbf 0$ 非负，故一阶条件成立。\emph{二阶（$\Rightarrow$）}：若某点存在方向 $\mathbf v$ 使 $\mathbf v^\top\nabla^2 f(\mathbf x_0)\mathbf v<0$，令 $g(t)=f(\mathbf x_0+t\mathbf v)$，则 $g''(0)<0$，$g$ 在该点附近严格凹，与 “凸函数限制在直线上仍凸” 矛盾。
\end{proof}

\begin{pitfall}[思维陷阱：二阶条件不是判凸的唯一方法]
二阶条件仅适用于二次可微函数。$f(\mathbf x)=\|\mathbf x\|_1$ 是凸函数，但在 $\mathbf x=\mathbf 0$ 不可微，无法用 \cref{eq:cvx-second-order}。对非光滑函数，须回到定义式 \cref{eq:convex-fun-def} 或 epigraph。另一个常见误区：“下水平集 $\{\mathbf x:f(\mathbf x)\le\alpha\}$ 凸” \emph{不能}反推 $f$ 凸——$f(\mathbf x)=\sqrt{\|\mathbf x\|}$ 的每个下水平集都是球（凸），但 $f$ 不凸（它只是\emph{拟凸}）。
\end{pitfall}

\subsection{凸函数的保凸运算}
\label{subsec:cvx-fun-ops}

实际目标函数往往由复杂运算组合而成。保凸运算让我们 “模块化” 地判定凸性——这正是 DCP 规则的数学基础。

\begin{theorem}[凸函数的保凸运算]\label{thm:fun-ops}
下列运算保凸：
\begin{enumerate}
  \item \textbf{非负加权和}：$f_i$ 凸、$w_i\ge0\Rightarrow\sum_i w_i f_i$ 凸；
  \item \textbf{仿射复合}：$f$ 凸 $\Rightarrow g(\mathbf x)=f(\mathbf A\mathbf x+\mathbf b)$ 凸；
  \item \textbf{逐点上确界}：$f_\alpha$ 皆凸 $\Rightarrow\sup_\alpha f_\alpha$ 凸；
  \item \textbf{部分最小化}：$f(\mathbf x,\mathbf y)$ 关于 $(\mathbf x,\mathbf y)$ 联合凸、$C$ 凸 $\Rightarrow g(\mathbf x)=\inf_{\mathbf y\in C}f(\mathbf x,\mathbf y)$ 凸（若 $>-\infty$）；
  \item \textbf{透视}：$f$ 凸 $\Rightarrow g(\mathbf x,t)=t\,f(\mathbf x/t)$（$t>0$）凸；
  \item \textbf{单调复合}：$h$ 凸且不减、$g$ 凸 $\Rightarrow h\circ g$ 凸（或 $h$ 凸且不增、$g$ 凹）。
\end{enumerate}
\end{theorem}
\begin{proof}
仅证逐点上确界（最具代表性）：$\operatorname{epi}(\sup_\alpha f_\alpha)=\bigcap_\alpha\operatorname{epi}(f_\alpha)$，凸集的交仍凸（\cref{thm:set-ops}），故 $\sup_\alpha f_\alpha$ 凸。其余把定义式 \cref{eq:convex-fun-def} 代入直接验证。
\end{proof}

\begin{insight}[$\sup$ 保凸而 $\inf$ 一般不保]\label{ins:sup-inf}
$\sup$ 运算保凸，但 $\inf$ 一般不保（除非\emph{联合}凸，即部分最小化的条件）。这个不对称解释了一个深刻现象：“最坏情况”（max/sup，鲁棒优化）容易在凸框架内处理，而 “最好情况”（min/inf）需要额外结构。Moreau 包络 $M_f^\lambda(\mathbf v)=\inf_{\mathbf x}\{f(\mathbf x)+\tfrac1{2\lambda}\|\mathbf x-\mathbf v\|^2\}$ 正是部分最小化的实例，它把非光滑的 $f$ 变成光滑的 $M_f^\lambda$，是 proximal 算法的核心（\cref{sec:cvx-prox}）。
\end{insight}

\begin{pitfall}[概念误区：复合只要 $h,g$ 都凸就凸]
复合 $h(g(\mathbf x))$ 的凸性需要\emph{单调性}配合：$h$ 凸\emph{且不减} $+$ $g$ 凸，或 $h$ 凸\emph{且不增} $+$ $g$ 凹。反例：$h(t)=t^2$ 凸、$g(x)=-x^2+1$ 凸，但 $h(g(x))=(1-x^2)^2$ 非凸。此外，凸函数之\emph{积}、之\emph{差} 一般不凸（$f=x,g=-x$ 时 $fg=-x^2$ 凹）。
\end{pitfall}

\subsection{严格凸、强凸与闭凸函数}
\label{subsec:cvx-fun-types}

\begin{definition}[凸 / 严格凸 / 强凸 / 闭凸]\label{def:cvx-strength}
对 $\mathbf x\neq\mathbf y$、$\theta\in(0,1)$：
\begin{itemize}
  \item \textbf{严格凸}：\cref{eq:convex-fun-def} 取严格不等号 “$<$”；
  \item \textbf{$\mu$-强凸}（$\mu>0$）：$f-\tfrac{\mu}{2}\|\cdot\|^2$ 是凸函数；
  \item \textbf{下半连续（lsc）}：$\liminf_{\mathbf x\to\mathbf x_0}f(\mathbf x)\ge f(\mathbf x_0)$（函数可向上跳、不可向下跳），等价于 $\operatorname{epi}f$ 是闭集；
  \item \textbf{闭凸函数}：proper（不恒 $+\infty$、不取 $-\infty$）$+$ 凸 $+$ lsc。
\end{itemize}
层级关系：强凸 $\Rightarrow$ 严格凸 $\Rightarrow$ 凸，反向均不成立。
\end{definition}

\begin{insight}[为何强凸远比严格凸重要]\label{ins:strong-vs-strict}
$f(x)=x^4$ 严格凸但\emph{不}强凸（最小处曲率为零）。严格凸只保证最小值点\emph{唯一}；强凸还要求曲率有正的下界（“至少像二次函数那样弯”），从而\emph{额外}给出梯度下降的\emph{线性收敛}与定量误差界（\cref{sec:cvx-strong}）。这就是为什么收敛性分析里我们真正依赖的是强凸而非严格凸。
\end{insight}

为什么需要 “闭”？因为对偶理论的核心定理 $f^{\ast\ast}=f$（\cref{sec:cvx-conjugate}）只对闭凸函数成立——若函数不 “闭”，共轭运算会丢失信息。一个安心的事实：有限维空间中的凸函数在 $\operatorname{ri}(\operatorname{dom}f)$ 上\emph{自动连续}（凸性这个代数性质蕴含了连续性这个分析性质），故范数、二次函数、指示函数等常见函数都是闭凸的；只有在定义域\emph{边界}上凸函数才可能不连续（如指示函数 $\delta_C$ 在 $\partial C$ 上）。

\begin{practice}
\begin{enumerate}
  \item 用三种方法（定义、一阶、二阶）证明 $f(x)=e^x$ 凸。
  \item 证明 log-sum-exp $f(\mathbf x)=\log\sum_i e^{x_i}$ 凸。（提示：算 Hessian $\nabla^2 f=\operatorname{diag}(\mathbf p)-\mathbf p\mathbf p^\top$，$p_i=e^{x_i}/\sum_j e^{x_j}$，则 $\mathbf v^\top\nabla^2 f\,\mathbf v=\operatorname{Var}_{\mathbf p}[\mathbf v]\ge0$。）
  \item 给出一个严格凸但不强凸的函数并证明其性质。
\end{enumerate}
\end{practice}

% ════════════════════════════════════════════════════════════════
\section{次梯度与最优性条件}
\label{sec:cvx-subgrad}

\paragraph{动机：把微分推广到非光滑凸函数。}
梯度是光滑优化的基本工具，但机器人优化中充满非光滑函数：$\ell_1$ 范数（稀疏优化）、指示函数（硬约束）、$\max$（minimax）、hinge loss（SVM）。这些函数在某些点不可微，梯度不存在。次梯度把微分工具推广到非光滑凸函数，是 ADMM、proximal gradient 等算法的理论基础。

\begin{definition}[次梯度与次微分]\label{def:subgradient}
设 $f$ 凸。向量 $\mathbf g$ 是 $f$ 在 $\mathbf x$ 处的\textbf{次梯度}，如果
\begin{equation}
  f(\mathbf y)\ge f(\mathbf x)+\mathbf g^\top(\mathbf y-\mathbf x),\qquad\forall\mathbf y\in\operatorname{dom}f.
  \label{eq:subgradient-def}
\end{equation}
所有次梯度之集合 $\partial f(\mathbf x)$ 称\textbf{次微分}。
\end{definition}

\noindent\textbf{几何直觉.}\;次梯度定义的仿射函数 $h(\mathbf y)=f(\mathbf x)+\mathbf g^\top(\mathbf y-\mathbf x)$ 是 $f$ 在点 $(\mathbf x,f(\mathbf x))$ 处的一个\emph{支撑超平面}——它从下方 “托住” 整个函数图形（\cref{fig:cvx-subgrad}）。次梯度之于凸函数，就像切线之于光滑曲线：光滑曲线每点只有一条切线（梯度唯一），凸函数在 “棱” 上可以有多条支撑线（次梯度不唯一）。若 $f$ 在 $\mathbf x$ 可微，则 $\partial f(\mathbf x)=\{\nabla f(\mathbf x)\}$ 退化为单点。

\begin{figure}[ht]
\centering
\begin{tikzpicture}[scale=1.05,>=Stealth]
  \draw[->] (-2.4,0)--(2.6,0) node[right, font=\scriptsize] {$x$};
  \draw[->] (0,-1.3)--(0,2.4) node[above, font=\scriptsize] {$f$};
  % 过原点的支撑线族：斜率 g in [-1,1]，每条都在 |x| 下方
  \foreach \g in {-1,-0.5,0,0.5,1}{
    \draw[second, thick, dashed] (-2.1,{-2.1*\g})--(2.1,{2.1*\g});
  }
  % f(x) = |x|
  \draw[very thick, main] (-2.2,2.2)--(0,0)--(2.2,2.2);
  \node[main, font=\small] at (1.55,2.1) {$f(x)=|x|$};
  \node[second!60!black, font=\scriptsize] at (1.85,-1.0) {斜率 $g\in[-1,1]$};
  \node[second!60!black, font=\scriptsize] at (-1.7,1.55) {$\partial f(0)=[-1,1]$};
  \fill (0,0) circle (1.2pt);
\end{tikzpicture}
\caption{$f(x)=|x|$ 在 $x=0$ 处的次微分是整段斜率区间 $[-1,1]$：每条斜率 $g\in[-1,1]$ 的过原点直线（虚线族）都从下方支撑 $|x|$。}
\label{fig:cvx-subgrad}
\end{figure}

\subsection{典型次微分与计算规则}
\label{subsec:cvx-subgrad-calc}

\begin{table}[H]\centering\small
\caption{常见函数的次微分}
\label{tab:subdiff}
\begin{tabular}{lll}
\toprule
函数 & 次微分 & 备注 \\
\midrule
$|x|$ & $\{1\}\,(x>0)$；$\{-1\}\,(x<0)$；$[-1,1]\,(x=0)$ & 标量原型 \\
$\|\mathbf x\|_1$ & 各分量 $\operatorname{sign}$，$0$ 处取 $[-1,1]$ & 逐分量笛卡尔积 \\
$\delta_C(\mathbf x)$ & $N_C(\mathbf x)=\{\mathbf g:\mathbf g^\top(\mathbf y-\mathbf x)\le0,\forall\mathbf y\in C\}$ & 法锥 \\
$\max_i x_i$ & $\operatorname{conv}\{\mathbf e_i:i\in I(\mathbf x)\}$ & $I(\mathbf x)$ 为取最大下标集 \\
\bottomrule
\end{tabular}
\end{table}

次微分的计算规则让我们从简单函数组装复杂函数的次微分：
\begin{itemize}
  \item \textbf{求和}：$\partial(f_1+f_2)=\partial f_1+\partial f_2$（Minkowski 和），需正则条件 $\operatorname{ri}(\operatorname{dom}f_1)\cap\operatorname{ri}(\operatorname{dom}f_2)\neq\emptyset$；
  \item \textbf{数乘}：$\partial(\alpha f)=\alpha\partial f$（$\alpha>0$）；
  \item \textbf{仿射复合}：$\partial(f(\mathbf A\cdot+\mathbf b))(\mathbf x)=\mathbf A^\top\partial f(\mathbf A\mathbf x+\mathbf b)$（正则条件下）；
  \item \textbf{逐点最大}：$\partial(\max_i f_i)(\mathbf x)=\operatorname{conv}\bigcup_{i\in I(\mathbf x)}\partial f_i(\mathbf x)$，$I(\mathbf x)$ 为 active 集。
\end{itemize}

\begin{proposition}[次微分是闭凸集且单调]\label{prop:subdiff-monotone}
对凸 $f$，每点 $\partial f(\mathbf x)$ 是闭凸集；且 $\partial f$ 是\emph{单调算子}：对任意 $\mathbf g_x\in\partial f(\mathbf x)$、$\mathbf g_y\in\partial f(\mathbf y)$，
\begin{equation}
  \langle\mathbf g_x-\mathbf g_y,\ \mathbf x-\mathbf y\rangle\ge0.
  \label{eq:subdiff-monotone}
\end{equation}
\end{proposition}
\begin{proof}
闭凸：$\partial f(\mathbf x)=\bigcap_{\mathbf y}\{\mathbf g:\mathbf g^\top(\mathbf y-\mathbf x)\le f(\mathbf y)-f(\mathbf x)\}$ 是半空间的交。单调：把 \cref{eq:subgradient-def} 对 $(\mathbf x,\mathbf y)$ 与 $(\mathbf y,\mathbf x)$ 各写一遍相加即得。
\end{proof}

单调性是 proximal 算子理论的基石——正因 $\partial f$ 单调，$\operatorname{prox}_f=(\mathbf I+\partial f)^{-1}$ 才是良定义且非扩张的（\cref{sec:cvx-prox}）。

\subsection{Fermat 规则：一切最优性条件的根源}
\label{subsec:cvx-fermat}

\begin{theorem}[Fermat 规则]\label{thm:fermat}
$\mathbf x^\star$ 是凸函数 $f$ 的全局最小值点 $\iff\mathbf 0\in\partial f(\mathbf x^\star)$。
\end{theorem}
\begin{proof}
$\mathbf 0\in\partial f(\mathbf x^\star)$ 按 \cref{eq:subgradient-def} 即 $f(\mathbf y)\ge f(\mathbf x^\star)+\mathbf 0^\top(\mathbf y-\mathbf x^\star)=f(\mathbf x^\star)$ 对所有 $\mathbf y$ 成立，正是 “$\mathbf x^\star$ 全局最小”。两个方向互为定义改写。
\end{proof}

\begin{insight}[Fermat 规则统一了所有最优性条件]\label{ins:fermat-unify}
对凸函数，\emph{局部最优 $=$ 全局最优 $=$ $\mathbf 0\in\partial f(\mathbf x^\star)$} 三者等价，凸优化在理论上 “完全可解”。更关键的是它的\emph{普适性}：把约束 $\mathbf x\in C$ 写成指示函数 $\delta_C$，约束问题 $\min_{\mathbf x\in C}f$ 变成 $\min f+\delta_C$，Fermat 规则给出
\[
  \mathbf 0\in\nabla f(\mathbf x^\star)+N_C(\mathbf x^\star)\quad\Longleftrightarrow\quad-\nabla f(\mathbf x^\star)\in N_C(\mathbf x^\star),
\]
即 “负梯度落在可行域法锥内”。把这套用到标准约束结构 $f_i\le0,h_j=0$ 上展开，就\emph{自然长出} KKT 条件（\cref{sec:cvx-kkt}）——KKT 不是被 “发明” 的，而是 Fermat 规则在约束结构下的必然结果。
\end{insight}

\paragraph{应用：$\ell_1$ 正则为何稀疏。}
考虑 $\min F(\mathbf x)=\tfrac12\|\mathbf A\mathbf x-\mathbf b\|^2+\lambda\|\mathbf x\|_1$。Fermat 规则给 $\mathbf 0\in\mathbf A^\top(\mathbf A\mathbf x^\star-\mathbf b)+\lambda\partial\|\mathbf x^\star\|_1$。记残差梯度 $\mathbf r=\mathbf A^\top(\mathbf A\mathbf x^\star-\mathbf b)$，逐分量看：若 $|r_i|<\lambda$，则 $-r_i/\lambda\in(-1,1)$ 落在 $\partial|x_i|$ 在 $0$ 处的区间内部，唯一相容解是 $x_i^\star=0$。于是 $\lambda$ 越大，越多分量被 “阈值” 杀成零——这就是 $\ell_1$ 正则化产生稀疏解的数学原因，也预告了 \cref{sec:cvx-prox} 的软阈值算子。

\begin{pitfall}[思维陷阱：把 $\mathbf 0\in\partial f$ 读成 “导数为零”]
对光滑函数确为 $\nabla f(\mathbf x^\star)=\mathbf 0$；但对非光滑函数，$\partial f(\mathbf x^\star)$ 是一个\emph{集合}，$\mathbf 0\in\partial f(\mathbf x^\star)$ 表示\emph{零向量属于次微分集}，不是 “导数为零”。在不可微点 $\partial f$ 含多个向量，只要其中一个是零就够了。另外，数值差分 $(f(\mathbf x+h)-f(\mathbf x-h))/(2h)$ 在不可微点给出的是某方向的方向导数，\emph{不是}次梯度集合——非光滑处应按结构解析计算。
\end{pitfall}

\begin{practice}
\begin{enumerate}
  \item 用 Fermat 规则求 $\min\|\mathbf x\|_1+\tfrac12\|\mathbf x-\mathbf v\|^2$（即 proximal 算子），推出软阈值公式。
  \item 证明对可微凸 $f$，$\mathbf x^\star=\arg\min_{\mathbf x\in C}f$ 当且仅当 $\langle\nabla f(\mathbf x^\star),\mathbf y-\mathbf x^\star\rangle\ge0$（$\forall\mathbf y\in C$）。
  \item 用次微分条件给出 LASSO 中 “第 $i$ 分量为零” 的充要条件。
\end{enumerate}
\end{practice}

% ════════════════════════════════════════════════════════════════
\section{强凸性、光滑性与条件数}
\label{sec:cvx-strong}

\paragraph{动机：凸性管 “能否”，强凸与光滑管 “多快”。}
凸性保证 “局部最优 $=$ 全局最优”，但没说算法要多快才能找到它。\textbf{强凸性}（曲率下界）与\textbf{光滑性}（曲率上界）提供了定量回答——它们共同决定一阶算法的收敛速率。

\begin{definition}[$\mu$-强凸与 $L$-光滑]\label{def:strong-smooth}
对凸 $f$：
\begin{itemize}
  \item \textbf{$\mu$-强凸}（三等价，可微时）：$f-\tfrac\mu2\|\cdot\|^2$ 凸 $\iff f(\mathbf y)\ge f(\mathbf x)+\nabla f(\mathbf x)^\top(\mathbf y-\mathbf x)+\tfrac\mu2\|\mathbf y-\mathbf x\|^2\iff\nabla^2 f\succeq\mu\mathbf I$；
  \item \textbf{$L$-光滑}（三等价）：$\|\nabla f(\mathbf x)-\nabla f(\mathbf y)\|\le L\|\mathbf x-\mathbf y\|\iff f(\mathbf y)\le f(\mathbf x)+\nabla f(\mathbf x)^\top(\mathbf y-\mathbf x)+\tfrac L2\|\mathbf y-\mathbf x\|^2\iff\nabla^2 f\preceq L\mathbf I$。
\end{itemize}
\textbf{条件数} $\kappa=L/\mu\ge1$。
\end{definition}

\noindent\textbf{“二次夹逼” 直觉.}\;$\mu$-强凸 $+$ $L$-光滑意味着函数被两个二次函数夹住：$\tfrac\mu2\|\mathbf x-\mathbf x^\star\|^2\le f(\mathbf x)-f(\mathbf x^\star)\le\tfrac L2\|\mathbf x-\mathbf x^\star\|^2$。把凸函数比作一个碗，$\mu$ 是碗底最窄处的曲率（决定收敛多快），$L$ 是碗壁最陡处的曲率（决定步长不能太大，否则飞出碗外），$\kappa$ 衡量碗的 “扁平度”。

\begin{theorem}[下降引理 / Descent Lemma]\label{thm:descent-lemma}
若 $f$ 是 $L$-光滑的，则对任意 $\mathbf x,\mathbf y$，
\begin{equation}
  f(\mathbf y)\le f(\mathbf x)+\nabla f(\mathbf x)^\top(\mathbf y-\mathbf x)+\frac L2\|\mathbf y-\mathbf x\|^2.
  \label{eq:descent-lemma}
\end{equation}
\end{theorem}
\begin{proof}
由微积分基本定理 $f(\mathbf y)-f(\mathbf x)=\int_0^1\nabla f(\mathbf x+t(\mathbf y-\mathbf x))^\top(\mathbf y-\mathbf x)\,dt$。减去 $\nabla f(\mathbf x)^\top(\mathbf y-\mathbf x)$ 后，对积分项用 Cauchy--Schwarz 与 $L$-Lipschitz：$|\!\int_0^1[\nabla f(\mathbf x+t(\mathbf y-\mathbf x))-\nabla f(\mathbf x)]^\top(\mathbf y-\mathbf x)\,dt|\le\int_0^1 Lt\|\mathbf y-\mathbf x\|^2\,dt=\tfrac L2\|\mathbf y-\mathbf x\|^2$。
\end{proof}

下降引理是一阶方法的基石：取 $\mathbf y=\mathbf x-\tfrac1L\nabla f(\mathbf x)$（步长 $1/L$ 的梯度步）代入 \cref{eq:descent-lemma}，得 $f(\mathbf y)\le f(\mathbf x)-\tfrac1{2L}\|\nabla f(\mathbf x)\|^2$，即\emph{每步至少下降 $\tfrac1{2L}\|\nabla f(\mathbf x)\|^2$}。这就是步长 $\eta=1/L$ 的理论依据，\cref{sec:cvx-algo} 的全部收敛率证明都从这里出发。

\begin{theorem}[Co-coercivity / Baillon--Haddad]\label{thm:cocoercivity}
若 $f$ \emph{凸}且 $L$-光滑，则
\begin{equation}
  \langle\nabla f(\mathbf x)-\nabla f(\mathbf y),\ \mathbf x-\mathbf y\rangle\ge\frac1L\|\nabla f(\mathbf x)-\nabla f(\mathbf y)\|^2.
  \label{eq:cocoercivity}
\end{equation}
\end{theorem}

\begin{insight}[条件数决定一阶算法的迭代数]\label{ins:condition-number}
条件数 $\kappa=L/\mu$ 几乎单手决定一阶算法的速度。对 $\mu$-强凸 $L$-光滑问题，梯度下降的误差以 $(1-1/\kappa)$ 线性衰减，达 $\epsilon$ 精度需 $O(\kappa\log(1/\epsilon))$ 步；Nesterov 加速把 $\kappa$ 换成 $\sqrt\kappa$，只需 $O(\sqrt\kappa\log(1/\epsilon))$ 步。当 $\kappa=10^4$（机器人逆动力学 Hessian 的典型值）时，二者相差 $100$ 倍——这就是加速法的出场动机（\cref{sec:cvx-algo}）。\cref{thm:cocoercivity} 的 co-coercivity 则是证明近端梯度、forward--backward 分裂收敛的关键引理：它说梯度算子不仅单调，而且 “足够单调”（有定量下界）。
\end{insight}

\begin{pitfall}[编程陷阱：co-coercivity 需要凸性前提]
仅 $L$-光滑（梯度 Lipschitz）\textbf{不足}以保证 co-coercivity \cref{eq:cocoercivity}——非凸 $L$-光滑函数只有弱单调性 $\langle\nabla f(\mathbf x)-\nabla f(\mathbf y),\mathbf x-\mathbf y\rangle\ge-L\|\mathbf x-\mathbf y\|^2$。在非凸优化的收敛分析中不能直接套用凸情形结论。另一处常见错误：把步长写死 $\eta=0.01$ 而不估 $L$（$L=\lambda_{\max}(\nabla^2 f)$，对二次型即 $\lambda_{\max}(\mathbf Q)$），导致收敛极慢或振荡发散；实践中可用 Armijo 回溯自适应选步长。
\end{pitfall}

\begin{practice}
\begin{enumerate}
  \item 计算 $f(\mathbf x)=\tfrac12\mathbf x^\top\mathbf A\mathbf x+\mathbf b^\top\mathbf x$（$\mathbf A\succ\mathbf 0$）的 $\mu,L,\kappa$。
  \item 证明：$\mu$-强凸下梯度下降 $\mathbf x_{k+1}=\mathbf x_k-\tfrac1L\nabla f(\mathbf x_k)$ 满足 $\|\mathbf x_k-\mathbf x^\star\|^2\le(1-\mu/L)^k\|\mathbf x_0-\mathbf x^\star\|^2$。
  \item 对 $f(\mathbf x)=\tfrac12\mathbf x^\top\mathbf A\mathbf x$ 直接验证 co-coercivity 并找出 $L$。
\end{enumerate}
\end{practice}

% ════════════════════════════════════════════════════════════════
\section{凸优化问题与其家族}
\label{sec:cvx-problem}

\paragraph{动机：识别问题类型 $=$ 选对求解器。}
SLAM 后端用高斯--牛顿解非线性最小二乘、MPC 解 QP、certifiable perception 解 SDP——这些看似不同的问题有共同结构：它们都是（或可松弛为）凸优化问题。识别问题的 “类型” 是选择求解器、理解难度的第一步。

\begin{definition}[凸优化问题标准形式]\label{def:cvx-problem}
\begin{equation}
  \min_{\mathbf x}\ f_0(\mathbf x)\quad\text{s.t.}\quad f_i(\mathbf x)\le0\ (i=1,\dots,m),\quad h_j(\mathbf x)=0\ (j=1,\dots,p),
  \label{eq:cvx-standard}
\end{equation}
其中 $f_0,\dots,f_m$ 凸，$h_1,\dots,h_p$ \emph{仿射}。
\end{definition}

为什么等式约束必须仿射？若允许非线性等式 $h(\mathbf x)=0$，可行域一般不凸（如 $x_1^2+x_2^2=1$ 是圆，非凸集）；只有仿射等式 $\mathbf A\mathbf x=\mathbf b$（超平面的交）才保证凸性。

\begin{insight}[凸优化的根本保证：局部即全局]\label{ins:local-global}
凸优化问题的核心担保是：\emph{任何局部最优解都是全局最优解}。这不是 “运气好”，而是凸性的数学必然——若存在更好的可行点，凸性保证连接局部最优与该点的路径上函数值单调不增，矛盾。正是这条担保，让 “把问题凸化”（凸松弛、凸近似、凸包络、对偶下界）成为机器人学反复使用的策略。
\end{insight}

\paragraph{等价变换：把 “看起来不凸” 化为标准凸。}
凸优化的一项核心技能是变换。最常用的是 epigraph 变换：$\min f_0(\mathbf x)\iff\min_{\mathbf x,t}t$ s.t.\ $f_0(\mathbf x)\le t$。于是不可微的 minimax $\min_{\mathbf x}\max_i f_i(\mathbf x)$ 变成标准形式 $\min_{\mathbf x,t}t$ s.t.\ $f_i(\mathbf x)\le t$（线性目标 $+$ 凸约束）。类似地，$\min\|\mathbf A\mathbf x-\mathbf b\|_\infty$ 与 $\min\|\mathbf A\mathbf x-\mathbf b\|_1$ 都可化为 LP。初学者常在原变量空间苦思而忽略换元。

\subsection{六大家族与锥包含链}
\label{subsec:cvx-conic}

凸优化问题有精细的层级结构，从最简单的 LP 到最复杂的 SDP，每级有专用高效求解器。

\begin{table}[H]\centering\small
\caption{凸优化问题家族与机器人应用}
\label{tab:cvx-families}
\begin{tabular}{lll}
\toprule
类型 & 标准形式（目标 / 关键约束） & 机器人应用 \\
\midrule
LP & $\mathbf c^\top\mathbf x$；$\mathbf A\mathbf x\preceq\mathbf b$ & 任务分配、$\ell_1$/$\ell_\infty$ 优化 \\
QP & $\tfrac12\mathbf x^\top\mathbf P\mathbf x+\mathbf q^\top\mathbf x$（$\mathbf P\succeq\mathbf 0$） & MPC、CBF-QP、WBC、逆运动学 \\
SOCP & $\mathbf f^\top\mathbf x$；$\|\mathbf A_i\mathbf x+\mathbf b_i\|\le\mathbf c_i^\top\mathbf x+d_i$ & 摩擦锥、鲁棒约束、火箭着陆 \\
SDP & $\langle\mathbf C,\mathbf X\rangle$；$\mathbf X\succeq\mathbf 0$ & SE-Sync、LMI 控制综合、SOS \\
\bottomrule
\end{tabular}
\end{table}

\begin{insight}[锥包含链 LP $\subset$ QP $\subset$ SOCP $\subset$ SDP]\label{ins:cone-chain}
这四类问题严格逐级包含：$\mathrm{LP}\subset\mathrm{QP}\subset\mathrm{QCQP}\subset\mathrm{SOCP}\subset\mathrm{SDP}$，复杂度也逐级上升（LP 约 $O(n^{2.5}L)$，SDP 最坏 $O(n^6)$）。它们都可统一到\emph{锥规划} $\min\langle\mathbf c,\mathbf x\rangle$ s.t.\ $\mathbf A\mathbf x=\mathbf b,\mathbf x\in K$，其中 $K$ 分别取 $\mathbb{R}^n_+$、$\mathrm{SOC}$ 积、$\mathbb{S}^n_+$。由于这些锥\emph{自对偶}（\cref{subsec:cvx-sep-cor}），锥规划的对偶 $\max\mathbf b^\top\mathbf y$ s.t.\ $\mathbf c-\mathbf A^\top\mathbf y\in K^\ast=K$ 与原问题同构——这是 Fenchel--Moreau 在锥规划中的体现。Schur 互补是连接各级的桥：它把 SOC 约束写成 LMI，把 QP 嵌进 SDP。
\end{insight}

\paragraph{凸松弛实例：Shor 松弛与 SE-Sync。}
非凸的二次约束二次规划（QCQP）$\min\mathbf x^\top\mathbf P_0\mathbf x+\mathbf q_0^\top\mathbf x$ s.t.\ $\mathbf x^\top\mathbf P_i\mathbf x+\mathbf q_i^\top\mathbf x\le b_i$，当某些 $\mathbf P_i$ 非 PSD 时是 NP-hard。令 $\mathbf X=\mathbf x\mathbf x^\top$，把 $\mathbf x^\top\mathbf P_i\mathbf x=\operatorname{tr}(\mathbf P_i\mathbf X)$，再把 $\mathbf X=\mathbf x\mathbf x^\top$ 松弛为 $\mathbf X\succeq\mathbf 0$（去掉 rank-1 约束），即得 \textbf{Shor 松弛}（一个 SDP）。松弛 “紧”（tight）的判据是 $\operatorname{rank}(\mathbf X^\star)=1$，此时从 $\mathbf X^\star$ 提取的 $\mathbf x^\star$ 就是全局最优。SLAM 中旋转同步 $\min\sum\|\mathbf R_i-\mathbf R_{ij}\mathbf R_j\|^2$ s.t.\ $\mathbf R^\top\mathbf R=\mathbf I$ 正是非凸 QCQP，Shor 松弛为 SDP 即 SE-Sync——它给出旋转估计的\emph{全局最优证书}。

\begin{pitfall}[思维陷阱：把 LP 的强对偶直觉移植到 SDP]
LP 只要原问题可行且有界就有强对偶；SDP \emph{可以}有有限的正对偶间隙（Ramana 的经典例子）。原因是半正定锥的相对内部可能与约束仿射空间只交于边界。SDP 的强对偶需要 Slater 条件（严格可行解存在）——见 \cref{sec:cvx-duality}。
\end{pitfall}

\begin{practice}
\begin{enumerate}
  \item 把 $\min\|\mathbf A\mathbf x-\mathbf b\|_\infty$ 与 $\min\|\mathbf A\mathbf x-\mathbf b\|_1$ 各化为标准 LP。
  \item 证明凸优化问题的最优解集是凸集；并证明若 $f_0$ 严格凸则最多一个最优解。
  \item 用 Schur 互补把 SOC 约束 $\|\mathbf A\mathbf x+\mathbf b\|\le\mathbf c^\top\mathbf x+d$ 写成 LMI，从而说明 SOCP $\subset$ SDP。
\end{enumerate}
\end{practice}

% ════════════════════════════════════════════════════════════════
\section{Lagrange 对偶与弱对偶}
\label{sec:cvx-duality}

\paragraph{动机：能否从外部给最优值一个下界？}
对偶理论的核心问题：能否从 “外部” 给出原问题最优值的下界？若能，且下界恰等于最优值（强对偶），我们就有了判定最优性的工具，以及非凸问题的全局最优\emph{证书}。

\begin{definition}[Lagrangian 与对偶函数]\label{def:lagrangian}
对问题 \cref{eq:cvx-standard}，\textbf{Lagrangian} 与\textbf{对偶函数}为
\begin{align}
  L(\mathbf x,\boldsymbol\lambda,\boldsymbol\nu)&=f_0(\mathbf x)+\sum_{i=1}^m\lambda_i f_i(\mathbf x)+\sum_{j=1}^p\nu_j h_j(\mathbf x),\label{eq:lagrangian}\\
  g(\boldsymbol\lambda,\boldsymbol\nu)&=\inf_{\mathbf x}L(\mathbf x,\boldsymbol\lambda,\boldsymbol\nu),\label{eq:dual-function}
\end{align}
其中 $\lambda_i\ge0$ 是不等式约束的对偶变量（Lagrange 乘子），$\nu_j$ 是等式约束的对偶变量。
\end{definition}

\noindent\textbf{直觉.}\;Lagrangian 通过 “软化” 约束获得下界——把硬约束 $f_i\le0$ 变成惩罚项 $\lambda_i f_i$：约束被违反（$f_i>0$）时惩罚为正、增大目标；约束满足（$f_i<0$）时 “奖励” 为负、减小目标。\emph{经济学解释}：$\lambda_i$ 是约束 $f_i\le0$ 的 “价格”（影子价格）——若放松约束（允许 $f_i\le\epsilon$），最优值大约改善 $\lambda_i\epsilon$。

\begin{theorem}[弱对偶]\label{thm:weak-duality}
对偶函数 $g$ 永远\emph{凹}（无论原问题是否凸）；且对任意 $\boldsymbol\lambda\succeq\mathbf 0$、$\boldsymbol\nu$，
\begin{equation}
  g(\boldsymbol\lambda,\boldsymbol\nu)\le p^\star,
  \label{eq:weak-duality}
\end{equation}
其中 $p^\star$ 是原问题最优值。
\end{theorem}
\begin{proof}
\emph{凹}：对每个固定 $\mathbf x$，$L(\mathbf x,\cdot,\cdot)$ 关于 $(\boldsymbol\lambda,\boldsymbol\nu)$ 仿射（既凸又凹），仿射函数族的 $\inf$ 保凹（对偶于 “$\sup$ 保凸”）。\emph{下界}：设 $\tilde{\mathbf x}$ 任意可行，则 $f_i(\tilde{\mathbf x})\le0,h_j(\tilde{\mathbf x})=0$，于是
\[
  g(\boldsymbol\lambda,\boldsymbol\nu)\le L(\tilde{\mathbf x},\boldsymbol\lambda,\boldsymbol\nu)=f_0(\tilde{\mathbf x})+\sum_i\underbrace{\lambda_i}_{\ge0}\underbrace{f_i(\tilde{\mathbf x})}_{\le0}+\sum_j\nu_j\underbrace{h_j(\tilde{\mathbf x})}_{=0}\le f_0(\tilde{\mathbf x}).
\]
对所有可行 $\tilde{\mathbf x}$ 取 inf 即得 $g\le p^\star$。
\end{proof}

\textbf{对偶问题}是在所有下界中找最紧的：$d^\star=\max_{\boldsymbol\lambda\succeq\mathbf 0,\boldsymbol\nu}g(\boldsymbol\lambda,\boldsymbol\nu)$。\textbf{对偶间隙} $p^\star-d^\star\ge0$（由弱对偶）；若间隙为零则称\emph{强对偶}成立。

\begin{insight}[对偶问题永远凸——即使原问题非凸]\label{ins:dual-always-convex}
由 \cref{thm:weak-duality}，$g$ 凹，故对偶问题 $\max g$ 永远是\emph{凸}优化——\textbf{即使原问题非凸}。这使得对偶下界在非凸问题中依然有价值：它是 certifiable 算法（如 SE-Sync）输出全局最优\emph{证书}的来源。写论文时把 “对偶给出下界” 误说成 “对偶等于最优” 是最频繁的错误——弱对偶 $g\le p^\star$ 对任何问题永远成立，而强对偶 $d^\star=p^\star$ 需要凸性 $+$ 约束规范（\cref{subsec:cvx-slater}）。
\end{insight}

\paragraph{速算示例.}\;$\min\|\mathbf x\|^2$ s.t.\ $\mathbf A\mathbf x=\mathbf b$。Lagrangian $L=\|\mathbf x\|^2+\boldsymbol\nu^\top(\mathbf A\mathbf x-\mathbf b)$，$\nabla_{\mathbf x}L=\mathbf 0$ 给 $\mathbf x=-\tfrac12\mathbf A^\top\boldsymbol\nu$，代回得无约束对偶 $\max_{\boldsymbol\nu}-\tfrac14\boldsymbol\nu^\top\mathbf A\mathbf A^\top\boldsymbol\nu-\mathbf b^\top\boldsymbol\nu$，解出 $\boldsymbol\nu^\star=-2(\mathbf A\mathbf A^\top)^{-1}\mathbf b$，于是 $\mathbf x^\star=\mathbf A^\top(\mathbf A\mathbf A^\top)^{-1}\mathbf b=\mathbf A^\dagger\mathbf b$——最小范数解\emph{不是凑出来的}，而是对偶最优条件的直接产物。

\subsection{Slater 条件与强对偶}
\label{subsec:cvx-slater}

\begin{definition}[Slater 条件]\label{def:slater}
凸问题 \cref{eq:cvx-standard} 满足 \textbf{Slater 条件}，如果存在严格可行点 $\tilde{\mathbf x}$：$f_i(\tilde{\mathbf x})<0$（所有非仿射的不等式约束严格满足）且 $\mathbf A\tilde{\mathbf x}=\mathbf b$。（对\emph{仿射}不等式约束只需 “$\le$” 即可——sharpened Slater。）
\end{definition}

\begin{theorem}[Slater $\Rightarrow$ 强对偶]\label{thm:strong-duality}
若凸问题满足 Slater 条件，则强对偶成立 $d^\star=p^\star$，且对偶最优被取到。
\end{theorem}

\begin{derivation}[强对偶的几何证明骨架]
定义值函数视角的集合 $\mathcal A=\{(\mathbf u,t):\exists\mathbf x,\ f_i(\mathbf x)\le u_i,\ h_j(\mathbf x)=0,\ f_0(\mathbf x)\le t\}$。由 $f_i$ 凸、$h_j$ 仿射，$\mathcal A$ 凸；且 $p^\star=\inf\{t:(\mathbf 0,t)\in\mathcal A\}$。令 $\mathcal B=\{(\mathbf 0,s):s<p^\star\}$，则 $\mathcal A\cap\mathcal B=\emptyset$。

对凸集 $\mathcal A,\mathcal B$ 用分离超平面定理（\cref{thm:separation}），存在非零 $(\tilde{\boldsymbol\lambda},\mu)$ 使 $\tilde{\boldsymbol\lambda}^\top\mathbf u+\mu t\ge\mu s$ 对所有 $(\mathbf u,t)\in\mathcal A$、$s<p^\star$ 成立，并可证 $\tilde{\boldsymbol\lambda}\succeq\mathbf 0$。\emph{关键一步是 $\mu>0$}：若 $\mu=0$，分离条件退化为 $\tilde{\boldsymbol\lambda}^\top\mathbf u\ge0$ 对所有可行松弛 $\mathbf u$ 成立；但 Slater 给出严格可行点 $\tilde{\mathbf x}$，其 $\tilde{\mathbf u}=(f_1(\tilde{\mathbf x}),\dots)\prec\mathbf 0$，于是 $\tilde{\boldsymbol\lambda}^\top\tilde{\mathbf u}\le0$，等号仅当 $\tilde{\boldsymbol\lambda}=\mathbf 0$——与 $(\tilde{\boldsymbol\lambda},\mu)\neq\mathbf 0$ 矛盾，故 $\mu>0$。令 $\boldsymbol\lambda=\tilde{\boldsymbol\lambda}/\mu$，分离条件给出 $g(\boldsymbol\lambda,\boldsymbol\nu)\ge p^\star$；结合弱对偶 $g\le p^\star$ 得 $g(\boldsymbol\lambda,\boldsymbol\nu)=p^\star$。
\end{derivation}

\begin{insight}[Slater 条件的几何含义：可行域有 “厚度”]\label{ins:slater-thickness}
Slater 条件的几何含义是可行域有 “厚度”（不是贴着约束边界的薄片）。这保证分离超平面不能 “水平”（$\mu=0$），从而目标函数值在分离中真正起作用，使对偶变量有限可达。换言之，把值函数 $p^\star(\mathbf u)=\inf\{f_0:f_i\le u_i\}$ 看作 $\mathbf u$ 的凸函数，强对偶 $d^\star=p^\star(\mathbf 0)$ 等价于 $p^\star$ 在 $\mathbf 0$ 处 “闭凸行为良好”——Slater 排除了 $p^\star$ 在 $\mathbf 0$ 处向下跳跃的退化。失去 Slater，整个 KKT 理论都可能失效。
\end{insight}

\begin{pitfall}[概念误区：混淆 Slater 与一般 NLP 的约束规范]
LICQ（active 约束梯度线性无关）、MFCQ 是\emph{一般 NLP}（Nocedal--Wright 语境）的约束规范；Slater 是\emph{凸问题专用}的，且最弱。对凸问题：只有不等式约束时 Slater $=$ 存在 $\bar{\mathbf x}$ 使 $f_i(\bar{\mathbf x})<0$；只有仿射约束时\emph{不需要} Slater，可行域非空即强对偶。不检查约束规范的后果：可能写出 KKT 却找不到满足它的点，或被求解器误报 “primal infeasible”，或得到错误的影子价格。
\end{pitfall}

\begin{practice}
\begin{enumerate}
  \item 对 $\min x^2$ s.t.\ $x\ge1$，写 Lagrangian、求 $g(\lambda)$、验证弱对偶与强对偶。
  \item 对 LP $\min\mathbf c^\top\mathbf x$ s.t.\ $\mathbf A\mathbf x=\mathbf b,\mathbf x\succeq\mathbf 0$ 推导其对偶。
  \item 解释为什么 “凸问题 $+$ 仿射约束” 不需要 Slater 即有强对偶。
\end{enumerate}
\end{practice}

% ════════════════════════════════════════════════════════════════
\section{KKT 条件}
\label{sec:cvx-kkt}

\paragraph{动机：约束优化的 “终极检验”。}
KKT 条件是约束优化的终极判据——满足它的点（在适当条件下）就是最优解。MPC 求解器（OSQP/qpOASES）的终止条件正是 KKT 残差足够小。

\begin{theorem}[KKT 条件]\label{thm:kkt}
设原问题凸且满足 Slater 条件。$\mathbf x^\star$ 是最优解、$(\boldsymbol\lambda^\star,\boldsymbol\nu^\star)$ 是对偶最优解，当且仅当下列四组条件同时成立：
\begin{subequations}\label{eq:kkt}
\begin{align}
  &\text{原始可行：}\quad f_i(\mathbf x^\star)\le0,\quad h_j(\mathbf x^\star)=0;\label{eq:kkt-primal}\\
  &\text{对偶可行：}\quad\lambda_i^\star\ge0;\label{eq:kkt-dual}\\
  &\text{互补松弛：}\quad\lambda_i^\star f_i(\mathbf x^\star)=0,\quad\forall i;\label{eq:kkt-comp}\\
  &\text{稳定性：}\quad\nabla f_0(\mathbf x^\star)+\sum_i\lambda_i^\star\nabla f_i(\mathbf x^\star)+\sum_j\nu_j^\star\nabla h_j(\mathbf x^\star)=\mathbf 0.\label{eq:kkt-stat}
\end{align}
\end{subequations}
（不可微时，稳定性即 $\mathbf 0\in\partial_{\mathbf x}L(\mathbf x^\star,\boldsymbol\lambda^\star,\boldsymbol\nu^\star)$。）
\end{theorem}

\begin{proof}
\emph{充分性}：设 KKT 成立。由稳定性 \cref{eq:kkt-stat} 与 $L(\cdot,\boldsymbol\lambda^\star,\boldsymbol\nu^\star)$ 关于 $\mathbf x$ 凸，$\mathbf x^\star$ 是 $L$ 的全局最小，故
\[
  p^\star\ge g(\boldsymbol\lambda^\star,\boldsymbol\nu^\star)=L(\mathbf x^\star,\boldsymbol\lambda^\star,\boldsymbol\nu^\star)=f_0(\mathbf x^\star)+\underbrace{\textstyle\sum_i\lambda_i^\star f_i(\mathbf x^\star)}_{=0\,(\text{互补})}+\underbrace{\textstyle\sum_j\nu_j^\star h_j(\mathbf x^\star)}_{=0}=f_0(\mathbf x^\star);
\]
又 $\mathbf x^\star$ 可行故 $f_0(\mathbf x^\star)\ge p^\star$，于是 $f_0(\mathbf x^\star)=p^\star$，$\mathbf x^\star$ 最优。

\emph{必要性}：由 Slater $\Rightarrow$ 强对偶（\cref{thm:strong-duality}），$g(\boldsymbol\lambda^\star,\boldsymbol\nu^\star)=p^\star=f_0(\mathbf x^\star)$。于是 $\inf_{\mathbf x}L(\mathbf x,\boldsymbol\lambda^\star,\boldsymbol\nu^\star)=f_0(\mathbf x^\star)$；而 $L(\mathbf x^\star,\boldsymbol\lambda^\star,\boldsymbol\nu^\star)=f_0(\mathbf x^\star)+\sum_i\lambda_i^\star f_i(\mathbf x^\star)\le f_0(\mathbf x^\star)$（因 $\lambda_i^\star\ge0,f_i\le0$）。两边夹逼得 $L(\mathbf x^\star,\dots)=f_0(\mathbf x^\star)$，故 $\sum_i\lambda_i^\star f_i(\mathbf x^\star)=0$；每项 $\le0$ 而和为零，逼出每项为零，即互补松弛 \cref{eq:kkt-comp}。又 $\mathbf x^\star$ 达到 $\inf_{\mathbf x}L$，由 Fermat 规则（\cref{thm:fermat}）得稳定性 \cref{eq:kkt-stat}。
\end{proof}

\noindent\textbf{各条件的直觉.}\;\emph{互补松弛}说：每个不等式约束，要么 active（$f_i=0$），要么其乘子为零（$\lambda_i^\star=0$）——经济上即 “资源没用尽就不值钱，只有用尽的资源才有影子价格”。\emph{稳定性}说：在最优解处，目标的负梯度可表示为 active 约束法向量的锥组合；若负梯度不在约束法锥中，就存在可行下降方向，$\mathbf x^\star$ 不最优。

\begin{theorem}[敏感性：对偶变量是影子价格]\label{thm:sensitivity}
设扰动问题最优值 $p^\star(\mathbf u)=\inf\{f_0(\mathbf x):f_i(\mathbf x)\le u_i,h_j(\mathbf x)=0\}$。强对偶下 $p^\star(\mathbf u)$ 是 $\mathbf u$ 的凸函数，且
\begin{equation}
  \lambda_i^\star=-\frac{\partial p^\star}{\partial u_i}\Big|_{\mathbf u=\mathbf 0},\qquad p^\star(\mathbf u)\ge p^\star-\boldsymbol\lambda^{\star\top}\mathbf u.
  \label{eq:sensitivity}
\end{equation}
\end{theorem}

\begin{insight}[对偶变量是天然的 “约束重要性排名”]\label{ins:shadow-price}
由 \cref{thm:sensitivity}，第 $i$ 个约束的对偶变量 $\lambda_i^\star$ 等于该约束放松一个单位时最优值的改善量。在 MPC-QP 里，约束的对偶变量直接告诉你 “哪个约束最紧、放松它能换多大性能”——这是调试与约束软化设计的利器：花时间去松一个 $\lambda_i^\star=0$ 的 inactive 约束毫无意义，应优先松 $\lambda_i^\star$ 大的那个。OSQP 返回的对偶变量 $\mathbf y$ 正是为此而生。
\end{insight}

\begin{pitfall}[概念误区：把 KKT 当万能充分条件]
KKT 作为\emph{充分}条件\textbf{只在凸问题中成立}；非凸问题的 KKT 点可能是局部极大、鞍点或伪最优。KKT 作为\emph{必要}条件则需约束规范（凸问题用 Slater；一般 NLP 用 LICQ/MFCQ）。CBF-QP 是凸的所以 KKT 充分，但一般非线性 MPC 的 SQP 内层只满足必要。另注意符号约定：不等式 $f_i\le0$ 配 $\lambda_i\ge0$，等式 $h_j=0$ 配无符号 $\nu_j$；不同求解器/教材可能定义 $L=f_0-\sum\lambda_i f_i$，用前务必对照文档。
\end{pitfall}

\begin{practice}
\begin{enumerate}
  \item 对 $\min x_1^2+x_2^2$ s.t.\ $x_1+x_2\ge1$，写 KKT 并求解。
  \item 对 OSQP 格式 QP $\min\tfrac12\mathbf x^\top\mathbf P\mathbf x+\mathbf q^\top\mathbf x$ s.t.\ $\mathbf l\preceq\mathbf A\mathbf x\preceq\mathbf u$，写完整 KKT 系统。
  \item 解释为何 OSQP 返回的对偶变量可用来判断 “哪个约束最值得放松”。
\end{enumerate}
\end{practice}

% ════════════════════════════════════════════════════════════════
\section{共轭函数与 Fenchel 对偶}
\label{sec:cvx-conjugate}

\paragraph{动机：把问题翻转到对偶视角。}
在凸优化中我们常需把问题 “翻转” 到等价视角——这就是对偶。\textbf{共轭函数（Legendre--Fenchel 变换）}是实现翻转的核心工具：它把 $f$ 映到 “对偶空间” 的 $f^\ast$，使很多困难问题（非光滑优化）在对偶空间变简单。对工程师而言这不是抽象玩具——OSQP/ADMM 每步迭代都涉及共轭运算，鲁棒 SLAM 的 GNC 本质是 Moreau 光滑化，RL 的 mirror descent 建立在 Fenchel 对偶上。

\begin{definition}[Legendre--Fenchel 共轭]\label{def:conjugate}
$f$ 的\textbf{共轭函数}为
\begin{equation}
  f^\ast(\mathbf y)=\sup_{\mathbf x\in\operatorname{dom}f}\big[\langle\mathbf y,\mathbf x\rangle-f(\mathbf x)\big].
  \label{eq:conjugate}
\end{equation}
\end{definition}

\noindent\textbf{几何意义：截距的上确界.}\;考虑斜率为 $\mathbf y$ 的仿射函数 $\ell(\mathbf x)=\langle\mathbf y,\mathbf x\rangle-b$，求最大的 $b$ 使 $\ell\le f$ 处处成立（$\ell$ 为 $f$ 的下方支撑线）。$\ell\le f\iff b\le\inf_{\mathbf x}[f(\mathbf x)-\langle\mathbf y,\mathbf x\rangle]=-f^\ast(\mathbf y)$。故 $f^\ast(\mathbf y)$ 是所有斜率 $\mathbf y$ 的支撑仿射函数中 “截距的负数的最大值”。共轭变换把函数从 “点值表示”（给 $\mathbf x$ 求 $f(\mathbf x)$）翻成 “切线表示”（给斜率 $\mathbf y$ 求最紧下方支撑线截距），就像 Fourier 变换把信号从时域翻到频域。

\begin{theorem}[Fenchel--Young 不等式]\label{thm:fenchel-young}
对任何 $f$ 与任何 $\mathbf x,\mathbf y$，
\begin{equation}
  f(\mathbf x)+f^\ast(\mathbf y)\ge\langle\mathbf x,\mathbf y\rangle,
  \label{eq:fenchel-young}
\end{equation}
等号成立当且仅当 $\mathbf y\in\partial f(\mathbf x)$。
\end{theorem}
\begin{proof}
由定义 $f^\ast(\mathbf y)\ge\langle\mathbf y,\mathbf x\rangle-f(\mathbf x)$，移项即得 \cref{eq:fenchel-young}。等号 $f^\ast(\mathbf y)=\langle\mathbf y,\mathbf x\rangle-f(\mathbf x)$ 表示 $\mathbf x$ 是 $\sup$ 的取到者，即 $f(\mathbf z)\ge f(\mathbf x)+\langle\mathbf y,\mathbf z-\mathbf x\rangle$ 对所有 $\mathbf z$ 成立，正是 $\mathbf y\in\partial f(\mathbf x)$。
\end{proof}

Fenchel--Young 不等式是所有对偶理论的出发点——弱对偶、Lagrange 对偶下界都可由它推出。

\begin{table}[H]\centering\small
\caption{常见函数的共轭}
\label{tab:conjugate}
\begin{tabular}{lll}
\toprule
$f(\mathbf x)$ & $f^\ast(\mathbf y)$ & 备注 \\
\midrule
$\tfrac12\mathbf x^\top\mathbf Q\mathbf x$（$\mathbf Q\succ\mathbf 0$） & $\tfrac12\mathbf y^\top\mathbf Q^{-1}\mathbf y$ & $\tfrac12\|\cdot\|^2$ 自共轭 \\
$\|\mathbf x\|_p$ & $\delta_{\{\|\mathbf y\|_q\le1\}}$（$\tfrac1p+\tfrac1q=1$） & 对偶范数球指示 \\
$\delta_C(\mathbf x)$ & $\sigma_C(\mathbf y)$ & 指示 $\to$ 支撑函数 \\
$-\log x$（$x>0$） & $-1-\log(-y)$（$y<0$） & \\
$\log\sum_i e^{x_i}$ & $\sum_i y_i\log y_i$（$\mathbf y\in\Delta_n$） & softmax 的共轭是负熵 \\
\bottomrule
\end{tabular}
\end{table}

\paragraph{共轭的基本性质.}\;$f^\ast$ \emph{永远凸且闭}（仿射函数族的逐点 $\sup$），即使 $f$ 非凸——这就是为什么对偶问题永远凸（\cref{ins:dual-always-convex}）。可分性 $(f_1\oplus f_2)^\ast=f_1^\ast+f_2^\ast$ 是 ADMM 分解的基础；infimal 卷积 $(f_1\square f_2)(\mathbf x)=\inf_{\mathbf z}[f_1(\mathbf z)+f_2(\mathbf x-\mathbf z)]$ 满足 $(f_1\square f_2)^\ast=f_1^\ast+f_2^\ast$，而 Moreau 包络恰是一个 infimal 卷积（\cref{sec:cvx-prox}）。

\begin{theorem}[Fenchel--Moreau 定理]\label{thm:fenchel-moreau}
对 proper 函数 $f$，
\begin{equation}
  f^{\ast\ast}=f\quad\Longleftrightarrow\quad f\ \text{是闭凸函数（proper $+$ lsc $+$ convex）}.
  \label{eq:fenchel-moreau}
\end{equation}
若 $f$ 非闭凸，则 $f^{\ast\ast}$ 等于 $f$ 的\emph{闭凸包}（最大的不超过 $f$ 的闭凸函数）。
\end{theorem}
\begin{proof}[证明骨架]
$f^{\ast\ast}\le f$ 总成立（$f^{\ast\ast}$ 是所有 $\le f$ 的仿射函数的逐点 $\sup$）。反向：若某点 $f^{\ast\ast}(\mathbf x_0)<f(\mathbf x_0)$，则 $(\mathbf x_0,f^{\ast\ast}(\mathbf x_0))$ 在闭凸集 $\operatorname{epi}f$ 的严格下方，由分离定理（\cref{thm:separation}）存在仿射 $\ell\le f$ 且 $\ell(\mathbf x_0)>f^{\ast\ast}(\mathbf x_0)$；但 $\ell\le f\Rightarrow\ell\le f^{\ast\ast}$，故 $\ell(\mathbf x_0)\le f^{\ast\ast}(\mathbf x_0)$，矛盾。
\end{proof}

\begin{insight}[闭凸世界在共轭下封闭]\label{ins:fenchel-moreau}
Fenchel--Moreau 定理说：\emph{闭凸函数的世界在共轭变换下封闭}——做两次共轭回到原函数。这使 “在原空间做优化” 与 “在对偶空间做优化” 完全等价。配套的共轭--次微分互逆定理 $\mathbf y\in\partial f(\mathbf x)\iff\mathbf x\in\partial f^\ast(\mathbf y)\iff f(\mathbf x)+f^\ast(\mathbf y)=\langle\mathbf x,\mathbf y\rangle$ 进一步给出：若 $f$ 可微严格凸，则 $\nabla f^\ast=(\nabla f)^{-1}$——共轭的梯度是原梯度的逆映射。这正是 mirror descent 与自然梯度的数学基础（信息几何中自然参数与均值参数的互换）。
\end{insight}

\subsection{强凸-光滑对偶}
\label{subsec:cvx-strong-smooth-dual}

\begin{theorem}[强凸--光滑对偶]\label{thm:strong-smooth-dual}
设 $f$ 闭凸。则
\begin{equation}
  f\ \text{是}\ \mu\text{-强凸}\quad\Longleftrightarrow\quad f^\ast\ \text{是}\ \tfrac1\mu\text{-光滑}.
  \label{eq:strong-smooth-dual}
\end{equation}
对称地，$f$ 是 $L$-光滑 $\iff f^\ast$ 是 $\tfrac1L$-强凸。
\end{theorem}
\begin{proof}[证明骨架（$\Rightarrow$）]
$f$ 是 $\mu$-强凸时 $\nabla f^\ast(\mathbf y)=\arg\min_{\mathbf x}[f(\mathbf x)-\langle\mathbf y,\mathbf x\rangle]$ 存在唯一。令 $\mathbf x_i=\nabla f^\ast(\mathbf y_i)$，则 $\mathbf y_i=\nabla f(\mathbf x_i)$。由 $f$ 的强单调性 $\langle\mathbf y_1-\mathbf y_2,\mathbf x_1-\mathbf x_2\rangle\ge\mu\|\mathbf x_1-\mathbf x_2\|^2$，配 Cauchy--Schwarz $\|\mathbf y_1-\mathbf y_2\|\,\|\mathbf x_1-\mathbf x_2\|\ge\mu\|\mathbf x_1-\mathbf x_2\|^2$，约去一个因子得 $\|\nabla f^\ast(\mathbf y_1)-\nabla f^\ast(\mathbf y_2)\|\le\tfrac1\mu\|\mathbf y_1-\mathbf y_2\|$，即 $\tfrac1\mu$-Lipschitz 梯度。
\end{proof}

\begin{insight}[条件数在共轭下不变]\label{ins:condition-invariant}
若 $f$ 既 $\mu$-强凸又 $L$-光滑（$\kappa=L/\mu$），则 $f^\ast$ 是 $\tfrac1L$-强凸且 $\tfrac1\mu$-光滑，条件数 $\kappa^\ast=\tfrac{1/\mu}{1/L}=L/\mu=\kappa$ \emph{不变}！这意味着 “在原空间做优化” 与 “在对偶空间做优化” 面临相同的困难度——所需迭代次数相同。这解释了为什么对偶加速方法不能靠 “切换到对偶空间” 改善条件数；真正的加速必须利用额外结构。一个反例式的妙用：给 LASSO 的 $\ell_1$ 项加 elastic-net 正则 $\tfrac\mu2\|\mathbf x\|^2$ 使原问题强凸，于是对偶问题变光滑、更易解。
\end{insight}

\begin{practice}
\begin{enumerate}
  \item 验证 $\big(\tfrac12\mathbf x^\top\mathbf Q\mathbf x\big)^\ast=\tfrac12\mathbf y^\top\mathbf Q^{-1}\mathbf y$（对 $\mathbf x$ 求导令零）。
  \item 推导 Huber 函数 $\rho_\delta$ 的共轭（答案 $\tfrac12 y^2+\delta_{[-\delta,\delta]}(y)$）。
  \item 验证 $f(\mathbf x)=\tfrac\mu2\|\mathbf x\|^2$ 是 $\mu$-强凸、其共轭 $\tfrac1{2\mu}\|\mathbf y\|^2$ 是 $\tfrac1\mu$-光滑。
\end{enumerate}
\end{practice}

% ════════════════════════════════════════════════════════════════
\section{Proximal 算子与 Moreau 包络}
\label{sec:cvx-prox}

\paragraph{动机：把对偶理论变成可执行迭代。}
共轭提供对偶视角，但要落地到算法需要一个能 “计算” 的工具。\textbf{Proximal 算子}就是把共轭--对偶理论变成可执行迭代的桥梁。光滑优化每步做显式梯度步 $\mathbf x_{k+1}=\mathbf x_k-\eta\nabla f(\mathbf x_k)$；若 $f$ 不可微（如 $\|\mathbf x\|_1$），梯度不存在——proximal 算子给出答案。

\begin{definition}[Proximal 算子与 Moreau 包络]\label{def:prox}
设 $f$ proper 闭凸。其\textbf{proximal 算子}与\textbf{Moreau 包络}为
\begin{align}
  \operatorname{prox}_{\lambda f}(\mathbf v)&=\arg\min_{\mathbf x}\Big[f(\mathbf x)+\frac1{2\lambda}\|\mathbf x-\mathbf v\|^2\Big],\label{eq:prox-def}\\
  M_f^\lambda(\mathbf v)&=\min_{\mathbf x}\Big[f(\mathbf x)+\frac1{2\lambda}\|\mathbf x-\mathbf v\|^2\Big]=\big(f\,\square\,\tfrac1{2\lambda}\|\cdot\|^2\big)(\mathbf v).\label{eq:moreau-def}
\end{align}
\end{definition}

\noindent\textbf{直觉.}\;proximal 算子找一个点，在 “靠近 $\mathbf v$”（被 $\tfrac1{2\lambda}\|\mathbf x-\mathbf v\|^2$ 惩罚）与 “使 $f$ 小” 之间折中。它像 “带弹簧的投影”——当 $f=\delta_C$ 时退化为正交投影 $\Pi_C$，一般情况加了软性的 $f$ 惩罚。把显式梯度步 $\mathbf x_{k+1}=\arg\min_{\mathbf x}[f(\mathbf x_k)+\nabla f(\mathbf x_k)^\top(\mathbf x-\mathbf x_k)+\tfrac1{2\eta}\|\mathbf x-\mathbf x_k\|^2]$（线性近似 $+$ 二次正则）换成隐式 proximal 步 $\mathbf x_{k+1}=\operatorname{prox}_{\eta f}(\mathbf x_k)$（直接用 $f$ 本身），就不再需要 $f$ 可微，数值上也更稳定。

\begin{theorem}[存在唯一性与 resolvent 刻画]\label{thm:prox-resolvent}
若 $f$ proper 闭凸，则 $\operatorname{prox}_f(\mathbf v)$ 对每个 $\mathbf v$ 存在且唯一，且
\begin{equation}
  \mathbf p=\operatorname{prox}_f(\mathbf v)\ \Longleftrightarrow\ \mathbf v-\mathbf p\in\partial f(\mathbf p)\ \Longleftrightarrow\ \operatorname{prox}_f=(\mathbf I+\partial f)^{-1}.
  \label{eq:prox-resolvent}
\end{equation}
\end{theorem}
\begin{proof}
\emph{存在唯一}：$g(\mathbf x)=f(\mathbf x)+\tfrac12\|\mathbf x-\mathbf v\|^2$ proper 闭且 $1$-强凸（二次项贡献 $\mu=1$），强凸 $\Rightarrow$ coercive $\Rightarrow\inf$ 取到，且严格凸 $\Rightarrow$ 唯一。\emph{resolvent}：由 Fermat 规则 $\mathbf 0\in\partial g(\mathbf p)=\partial f(\mathbf p)+(\mathbf p-\mathbf v)$，即 $\mathbf v-\mathbf p\in\partial f(\mathbf p)$，等价于 $\mathbf v\in(\mathbf I+\partial f)(\mathbf p)$，即 $\mathbf p=(\mathbf I+\partial f)^{-1}(\mathbf v)$。
\end{proof}

\begin{insight}[proximal 算子的四张面孔]\label{ins:prox-faces}
$\operatorname{prox}_f$ 是同一对象的四张面孔：\emph{优化}（解子问题）$=$ \emph{次微分}（残差是次梯度 $\mathbf v-\mathbf p\in\partial f(\mathbf p)$）$=$ \emph{resolvent}（单调算子的预解式 $(\mathbf I+\partial f)^{-1}$）$=$ \emph{不动点}（$\mathbf x^\star=\arg\min f\iff\mathbf x^\star=\operatorname{prox}_{\lambda f}(\mathbf x^\star)$）。最后这一面把最优性条件 $\mathbf 0\in\partial f$ 与不动点联系起来，是所有 proximal 算法收敛性证明的对应基础。配套的 \textbf{Moreau 分解} $\operatorname{prox}_f(\mathbf v)+\operatorname{prox}_{f^\ast}(\mathbf v)=\mathbf v$ 让我们由 $\operatorname{prox}_f$ 直接得 $\operatorname{prox}_{f^\ast}$，是 ADMM 在原--对偶之间往返的引擎。而由 $M_f^\lambda=f\square\tfrac1{2\lambda}\|\cdot\|^2$ 与 infimal 卷积的共轭规则，$(M_f^\lambda)^\ast=f^\ast+\tfrac\lambda2\|\cdot\|^2$——加二次项使共轭强凸，由强凸--光滑对偶（\cref{thm:strong-smooth-dual}），$M_f^\lambda$ \emph{总是光滑}的。这正是 Moreau 光滑化（鲁棒估计 GNC）的根源。
\end{insight}

\subsection{五个必记的闭式}
\label{subsec:cvx-prox-closed}

proximal 的实用性取决于能否高效计算。下面五个闭式覆盖绝大多数机器人优化场景。

\begin{table}[H]\centering\small
\caption{五个核心 proximal 算子（参数 $\lambda$）}
\label{tab:prox-closed}
\begin{tabular}{lll}
\toprule
$f(\mathbf x)$ & $\operatorname{prox}_{\lambda f}(\mathbf v)$ & 名称 / 用途 \\
\midrule
$\|\mathbf x\|_1$ & $\operatorname{sign}(\mathbf v)\odot(|\mathbf v|-\lambda)_+$ & 软阈值（LASSO，元素稀疏） \\
$\|\mathbf x\|_2$ & $\big(1-\lambda/\|\mathbf v\|\big)_+\mathbf v$ & 块软阈值（group LASSO，组稀疏） \\
$\delta_C(\mathbf x)$ & $\Pi_C(\mathbf v)$ & 投影（盒 / 球 / 单纯形） \\
$\tfrac12\mathbf x^\top\mathbf Q\mathbf x$ & $(\mathbf I+\lambda\mathbf Q)^{-1}\mathbf v$ & 二次（可预分解） \\
$\|\mathbf X\|_\ast$ & $\mathbf U\operatorname{diag}((\sigma_i-\lambda)_+)\mathbf W^\top$ & SVT（低秩补全、鲁棒 PCA） \\
\bottomrule
\end{tabular}
\end{table}

\begin{derivation}[软阈值：$\ell_1$ 的 prox]
$\|\cdot\|_1$ 可分，逐分量解 $p_i=\arg\min_{x_i}[\lambda|x_i|+\tfrac12(x_i-v_i)^2]$。由 Fermat 规则 $0\in\lambda\partial|x_i|+(x_i-v_i)$：
若 $x_i>0$ 则 $\partial|x_i|=\{1\}$，得 $x_i=v_i-\lambda$（需 $v_i>\lambda$）；若 $x_i<0$ 则 $x_i=v_i+\lambda$（需 $v_i<-\lambda$）；若 $x_i=0$ 则 $\partial|x_i|=[-1,1]$ 要求 $v_i\in[-\lambda,\lambda]$。合并即 $p_i=\operatorname{sign}(v_i)\,(|v_i|-\lambda)_+$。这就是 “软阈值”：绝对值小于 $\lambda$ 的分量被杀成零（产生稀疏），大于 $\lambda$ 的被收缩 $\lambda$——与 \cref{subsec:cvx-fermat} 的稀疏性分析完全吻合。
\end{derivation}

\begin{pitfall}[概念误区：范数 vs 范数平方的 prox 差异巨大]
$\operatorname{prox}_{\lambda\|\cdot\|_2}(\mathbf v)=(1-\lambda/\|\mathbf v\|)_+\mathbf v$（块软阈值，有死区，整组杀死或收缩）；而 $\operatorname{prox}_{\lambda\cdot\frac12\|\cdot\|^2}(\mathbf v)=\mathbf v/(1+\lambda)$（简单缩放，无死区）。混淆两者是 ADMM 实现中最常见的 bug。另注意 $\ell_1$ 是逐分量独立的软阈值，$\ell_2$ 是整组的块阈值——参数 $\lambda$ 的位置在不同论文/代码中也不统一，统一约定 $\operatorname{prox}_{\lambda f}$ 表示对 $\lambda f$ 取 prox、二次项系数 $\tfrac1{2\lambda}$。
\end{pitfall}

\begin{practice}
\begin{enumerate}
  \item 证明 $\operatorname{prox}_{\delta_C}(\mathbf v)=\Pi_C(\mathbf v)$、$\operatorname{prox}_{\frac12\|\cdot\|^2}(\mathbf v)=\mathbf v/2$。
  \item 用 Moreau 分解证明 $\operatorname{prox}_{\lambda\|\cdot\|_1}(\mathbf v)=\mathbf v-\lambda\,\Pi_{\{\|\cdot\|_\infty\le1\}}(\mathbf v/\lambda)$，其中 $\Pi_{\{\|\cdot\|_\infty\le1\}}$ 是到单位 $\ell_\infty$ 球的投影（即逐分量截断到 $[-1,1]$）。\footnote{此处对偶对象是 $\ell_1$ 范数的共轭 $(\|\cdot\|_1)^\ast=\delta_{\{\|\cdot\|_\infty\le1\}}$（对偶范数球的\emph{指示函数}，见 \cref{tab:conjugate}），而非 $\ell_\infty$ 范数本身；由 $\operatorname{prox}_{\delta_C}=\Pi_C$，缩放后的 Moreau 分解 $\operatorname{prox}_{\lambda f}(\mathbf v)+\lambda\operatorname{prox}_{\lambda^{-1}f^\ast}(\mathbf v/\lambda)=\mathbf v$ 把第二项化为到 $\ell_\infty$ 球的投影。}
  \item 推导块软阈值（提示：由旋转不变性，最优解在 $\mathbf v$ 方向上，化为标量问题）。
\end{enumerate}
\end{practice}

% ════════════════════════════════════════════════════════════════
\section{凸优化算法路线速览}
\label{sec:cvx-algo}

\paragraph{动机：算法选择 $=$ 识别结构 $+$ 查收敛率 $+$ 匹配预算。}
假设你为四足机器人设计 MPC，每 $20$ 毫秒要解一个 $200$ 变量、$500$ 约束的 QP。面前至少五种算法：梯度下降、加速梯度、ADMM、内点法、active-set。选哪个？选错了机器人要么求解超时摔倒、要么精度不够抖动。这不是 “试试看” 能回答的——每种算法都有严格收敛率，告诉我们在给定结构（光滑、强凸、可分、稀疏、规模）下到达 $\epsilon$ 精度需多少步、每步多贵。

\begin{insight}[结构换速度；黑箱一阶已被 Nesterov 封顶]\label{ins:structure-speed}
凸优化算法的核心叙事是 “\emph{结构换速度}”——你对问题结构知道得越多（光滑、强凸、可分、可 prox），就能选越快的算法。Nesterov 1983 年证明加速梯度达 $O(1/k^2)$ 并\emph{同时}证明这是黑箱一阶 oracle 下不可再改进的极限。此后所有 “加速” 都是在\emph{结构假设}（强凸、有限和、可 prox、稀疏）上换来的——黑箱世界已被封顶。不了解理论的工程师只能在黑箱求解器之间盲目切换。
\end{insight}

\begin{table}[H]\centering\small
\caption{凸优化主要算法的适用结构与收敛率}
\label{tab:cvx-algo}
\begin{tabular}{llll}
\toprule
算法 & 适用结构 & 收敛率（光滑凸 / 强凸） & 机器人落地 \\
\midrule
梯度下降 GD & 光滑无约束 & $O(1/k)$ / $(1-1/\kappa)^k$ & 一切的原子操作 \\
Nesterov 加速 AGD & 光滑（凸 / 强凸） & $O(1/k^2)$ / $(1-1/\sqrt\kappa)^k$ & 理论最优 \\
近端梯度 ISTA/FISTA & $f$ 光滑 $+$ $g$ 可 prox & $O(1/k)$ / $O(1/k^2)$ & LASSO、TV 去噪 \\
ADMM & 可分 $f(\mathbf x)+g(\mathbf z)$，各自可 prox & $O(1/k)$（常数极小，可分布） & OSQP（MPC-QP） \\
内点法 IPM & 中小规模、高精度、约束 & $O(\sqrt m\log(1/\epsilon))$ Newton 步 & HPIPM、acados（kHz MPC） \\
\bottomrule
\end{tabular}
\end{table}

\paragraph{四条主线一句话.}
\begin{itemize}
  \item \textbf{梯度下降}：$\mathbf x_{k+1}=\mathbf x_k-\tfrac1L\nabla f(\mathbf x_k)$，由下降引理（\cref{thm:descent-lemma}）每步保证下降；收敛率 $O(1/k)$（凸）、线性（强凸），完全由 $\kappa$ 决定。
  \item \textbf{Nesterov 加速}：先在动量方向\emph{外推} $\mathbf y_k=\mathbf x_k+\beta_k(\mathbf x_k-\mathbf x_{k-1})$ 再求梯度（lookahead），把 $O(1/k)$ 提升到 $O(1/k^2)$、把 $\kappa$ 换成 $\sqrt\kappa$。注意与 Polyak heavy-ball 的关键区别：heavy-ball 在一般凸函数上\emph{无} $O(1/k^2)$ 保证，AGD 在外推点求梯度且动量时变。
  \item \textbf{近端梯度 ISTA/FISTA}：对 $\min f+g$（$f$ 光滑、$g$ 可 prox），$\mathbf x_{k+1}=\operatorname{prox}_{\eta g}(\mathbf x_k-\eta\nabla f(\mathbf x_k))$；FISTA 外加 Nesterov 动量得 $O(1/k^2)$。LASSO 即 $f=\tfrac12\|\mathbf A\mathbf x-\mathbf b\|^2$、$g=\lambda\|\mathbf x\|_1$，prox 为软阈值。
  \item \textbf{ADMM}：对 $\min f(\mathbf x)+g(\mathbf z)$ s.t.\ $\mathbf A\mathbf x+\mathbf B\mathbf z=\mathbf c$，用增广 Lagrangian 交替更新 $\mathbf x$、$\mathbf z$、对偶 $\mathbf y$。OSQP 用它解 QP：KKT 矩阵在 MPC 闭环中不变，一次 LDL 分解 $+$ warm-start 使每次只需 $5$--$15$ 步。
  \item \textbf{内点法}：用 log-barrier 把约束问题转成一列无约束问题，每个用 Newton 法二次收敛求解；Newton 步数仅 $O(\log(1/\epsilon))$，是高精度 $+$ 约束场景的核心。
\end{itemize}

\begin{pitfall}[思维陷阱：混淆 “迭代复杂度” 与 “总计算时间”]
FISTA 只需 $O(\sqrt{L/\epsilon})$ 次迭代、IPM 需 $O(\sqrt m\log(1/\epsilon))$ 次 Newton 步，但\emph{总时间 $=$ 迭代数 $\times$ 每步代价}。FISTA 每步 $O(n)$（一次梯度 $+$ 一次 prox），IPM 每步 $O(n^3)$（一次 Newton 系统）。对 $n=100$ 的 QP，IPM 可能 $20$ 步结束，FISTA 可能要 $10^4$ 步。分水岭：\emph{中小规模高精度 $\to$ IPM；大规模中精度 $\to$ 一阶法}。同理 “越新的算法越好” 是错的——算法只有 “更适合”，先识别结构再查收敛率表。
\end{pitfall}

\begin{practice}
\begin{enumerate}
  \item 对 $f(\mathbf x)=\tfrac12(x_1^2+100x_2^2)$（$\kappa=100$），从 $\mathbf x_0=(1,1)$ 手推 GD 与 AGD 前 $5$ 步，对比轨迹。
  \item 解释为何共轭梯度在 $n$ 步内精确求解 $n$ 维二次规划而不违反 Nesterov 下界。
  \item 你需以 $50$ Hz 解 MPC-QP（$200$ 变量、$500$ 约束）。OSQP 与 qpOASES 哪个更合适？给出理由。
\end{enumerate}
\end{practice}

% ════════════════════════════════════════════════════════════════
\section*{本章小结}

\begin{table}[H]\centering\small
\caption{符号表}
\label{tab:cvx-symbols}
\begin{tabular}{lll}
\toprule
符号 & 含义 & 首次出现 \\
\midrule
$C,D$ & 凸集 & \cref{def:convex-set} \\
$\operatorname{conv},\operatorname{aff},\operatorname{ri}$ & 凸包 / 仿射包 / 相对内部 & \cref{def:convex-set,def:relint} \\
$K,K^\ast$ & 凸锥及其对偶锥（极锥） & \cref{ins:cone-hierarchy,subsec:cvx-sep-cor} \\
$\mathbb{S}^n_+,\mathrm{SOC}_n$ & 半正定锥 / 二阶锥 & \cref{prop:soc-psd-convex} \\
$\operatorname{epi}f,\operatorname{dom}f$ & 上图 / 有效定义域 & \cref{def:convex-fun} \\
$\delta_C,\sigma_C,N_C,\Pi_C$ & 指示 / 支撑函数 / 法锥 / 投影 & \cref{subsec:cvx-sep-cor,tab:subdiff} \\
$\partial f(\mathbf x)$ & 次微分 & \cref{def:subgradient} \\
$\mu,L,\kappa=L/\mu$ & 强凸参数 / 光滑常数 / 条件数 & \cref{def:strong-smooth} \\
$L(\mathbf x,\boldsymbol\lambda,\boldsymbol\nu),g(\boldsymbol\lambda,\boldsymbol\nu)$ & Lagrangian / 对偶函数 & \cref{def:lagrangian} \\
$p^\star,d^\star$ & 原 / 对偶最优值 & \cref{thm:weak-duality} \\
$f^\ast,f^{\ast\ast}$ & 共轭 / 双共轭 & \cref{def:conjugate,thm:fenchel-moreau} \\
$\operatorname{prox}_{\lambda f},M_f^\lambda$ & proximal 算子 / Moreau 包络 & \cref{def:prox} \\
\bottomrule
\end{tabular}
\end{table}

\begin{table}[H]\centering\small
\caption{定理速查表}
\label{tab:cvx-theorems}
\begin{tabular}{lll}
\toprule
定理 / 公式 & 一句话说明 & 对应 \\
\midrule
分离超平面 \cref{eq:separation} & 不相交凸集可被超平面分开 & \cref{thm:separation} \\
一阶 / 二阶判据 \cref{eq:cvx-first-order,eq:cvx-second-order} & 切线是全局下界；$\nabla^2 f\succeq\mathbf 0$ & \cref{thm:cvx-criteria} \\
保凸运算 & sup / 仿射 / 部分最小化保凸 & \cref{thm:fun-ops} \\
Fermat 规则 \cref{thm:fermat} & $\mathbf 0\in\partial f(\mathbf x^\star)\iff$ 全局最优 & \cref{thm:fermat} \\
下降引理 \cref{eq:descent-lemma} & 步长 $1/L$ 每步必降 & \cref{thm:descent-lemma} \\
co-coercivity \cref{eq:cocoercivity} & 梯度算子的定量单调性 & \cref{thm:cocoercivity} \\
弱对偶 \cref{eq:weak-duality} & 对偶给下界，对任何问题 & \cref{thm:weak-duality} \\
Slater $\Rightarrow$ 强对偶 & 严格可行 $\Rightarrow$ 零对偶间隙 & \cref{thm:strong-duality} \\
KKT 条件 \cref{eq:kkt} & 凸问题的最优性充要判据 & \cref{thm:kkt} \\
敏感性 \cref{eq:sensitivity} & 对偶变量 $=$ 影子价格 & \cref{thm:sensitivity} \\
Fenchel--Young \cref{eq:fenchel-young} & 对偶理论的出发点 & \cref{thm:fenchel-young} \\
Fenchel--Moreau \cref{eq:fenchel-moreau} & $f^{\ast\ast}=f\iff f$ 闭凸 & \cref{thm:fenchel-moreau} \\
强凸--光滑对偶 \cref{eq:strong-smooth-dual} & $\mu$-强凸 $\iff$ $\tfrac1\mu$-光滑 & \cref{thm:strong-smooth-dual} \\
prox resolvent \cref{eq:prox-resolvent} & $\operatorname{prox}_f=(\mathbf I+\partial f)^{-1}$ & \cref{thm:prox-resolvent} \\
\bottomrule
\end{tabular}
\end{table}

\begin{table}[H]\centering\small
\caption{知识点总表}
\label{tab:cvx-knowledge}
\begin{tabular}{clll}
\toprule
\# & 知识点 & 核心要点 & 难度 \\
\midrule
1 & 凸集 & 连线段在集合内；半空间是构件 & \(\bigstar\) \\
2 & 凸锥 / 锥包含链 & LP $\subset$ QP $\subset$ SOCP $\subset$ SDP & \(\bigstar\bigstar\) \\
3 & 分离超平面 & 对偶 / KKT 的几何根源 & \(\bigstar\bigstar\) \\
4 & 凸函数判据 & 一阶切线下界、二阶 Hessian $\succeq\mathbf 0$ & \(\bigstar\) \\
5 & 保凸运算 & sup / 仿射 / 部分最小化；DCP 基础 & \(\bigstar\bigstar\) \\
6 & 次梯度 & 非光滑广义导数；支撑超平面斜率 & \(\bigstar\bigstar\) \\
7 & Fermat 规则 & $\mathbf 0\in\partial f$；KKT 的统一根源 & \(\bigstar\bigstar\) \\
8 & 强凸 / 光滑 / 条件数 & 收敛率 $\sim\kappa$ 或 $\sqrt\kappa$ & \(\bigstar\bigstar\) \\
9 & 凸优化问题家族 & 识别类型 $=$ 选求解器 & \(\bigstar\bigstar\) \\
10 & Lagrange 对偶 / 弱对偶 & 对偶永远凸；下界与证书 & \(\bigstar\bigstar\) \\
11 & Slater / 强对偶 & 可行域有 “厚度”；$\mu>0$ & \(\bigstar\bigstar\bigstar\) \\
12 & KKT / 影子价格 & 四组件；敏感性分析 & \(\bigstar\bigstar\bigstar\) \\
13 & 共轭 / Fenchel--Moreau & 点值 $\leftrightarrow$ 切线表示；$f^{\ast\ast}=f$ & \(\bigstar\bigstar\bigstar\) \\
14 & 强凸--光滑对偶 & 条件数共轭下不变 & \(\bigstar\bigstar\bigstar\) \\
15 & proximal / Moreau 包络 & 四张面孔；五个闭式；光滑化 & \(\bigstar\bigstar\bigstar\) \\
16 & 算法路线 & 结构换速度；GD/FISTA/ADMM/IPM & \(\bigstar\bigstar\) \\
\bottomrule
\end{tabular}
\end{table}

% ════════════════════════════════════════════════════════════════
\section*{延伸阅读}
\begin{itemize}
  \item \textbf{工程化主线}：Boyd \& Vandenberghe《Convex Optimization》\cite{boyd2004convex}——凸集 / 凸函数 / 对偶 / 内点法的标准教材，本章凸优化问题家族、KKT、对偶几何证明的主要出处，附 CVXPY/DCP 的工程接口。
  \item \textbf{凸分析的奠基}：Rockafellar《Convex Analysis》\cite{rockafellar1970convex}——相对内部、共轭、次微分、回收锥的严格处理；低维凸集与 $\operatorname{ri}$ 的所有细节在此（Boyd 几乎不讲）。
  \item \textbf{一阶算法理论}：Nesterov《Lectures on Convex Optimization》\cite{nesterov2018lectures}——加速梯度、下界、estimate sequence 的权威推导；Beck《First-Order Methods in Optimization》\cite{beck2017fom}——ISTA/FISTA、proximal、对偶的统一教材。
  \item \textbf{proximal 与算子分裂}：Parikh \& Boyd《Proximal Algorithms》\cite{parikh2014proximal}（proximal 综述）、Bauschke \& Combettes《Convex Analysis and Monotone Operator Theory》\cite{bauschke2017convex}（单调算子、firmly nonexpansive、Moreau 分解的严格基础）。
  \item \textbf{ADMM 与分布式}：Boyd 等《Distributed Optimization via ADMM》\cite{boyd2011admm}——ADMM 推导、$\rho$ 自适应、Douglas--Rachford 等价，OSQP 的理论前身。
\end{itemize}

\section*{本章与后续章节的关系}
\begin{table}[H]\centering\small
\label{tab:cvx-downstream}
\begin{tabular}{lll}
\toprule
后续章节 & 关系 & 本章铺垫 \\
\midrule
\cref{ch:nlopt} 非线性优化 & 最小二乘、信赖域的凸性论证 & 凸函数判据、Fermat 规则、下降引理 \\
\cref{ch:slam_est} 状态估计 & MAP 估计 $=$ 加权最小二乘 & 凸性、对偶、$\ell_1$ 稀疏 \\
控制部 MPC / CBF-QP & QP 求解器（OSQP/qpOASES）原理 & QP 家族、KKT、ADMM、内点法 \\
控制部 H$\infty$ / LMI 综合 & Lyapunov 条件 $=$ SDP 可行性 & 半正定锥、SDP、强对偶 \\
certifiable SLAM（SE-Sync） & 旋转同步 SDP 松弛与最优证书 & Shor 松弛、对偶下界、自对偶锥 \\
\bottomrule
\end{tabular}
\end{table}

\section*{常见误解汇总}
\begin{table}[H]\centering\small
\label{tab:cvx-misconceptions}
\begin{tabular}{ll}
\toprule
误解 & 正确理解 \\
\midrule
“边界有角就不凸” & 凸性是代数性质（线性组合封闭），与边界光滑性无关 \\
“凸函数一定连续” & 仅在 $\operatorname{ri}(\operatorname{dom}f)$ 上自动连续，边界上可不连续 \\
“下水平集凸 $\Rightarrow$ 函数凸” & 只意味拟凸；$\sqrt{\|\mathbf x\|}$ 下水平集凸但非凸函数 \\
“凸函数之积仍凸” & 反例 $f=x,g=-x$，$fg=-x^2$ 凹 \\
“$\sup$ 与 $\inf$ 对凸性地位相同” & $\sup$ 保凸，$\inf$ 需联合凸（部分最小化） \\
“$\mathbf 0\in\partial f$ 表示导数为零” & 表示零向量\emph{属于}次微分集，非 “导数为零” \\
“弱对偶 $=$ 强对偶” & 弱对偶对任何问题成立；强对偶需凸 $+$ Slater \\
“KKT 是万能充分条件” & 仅凸问题充分；非凸 KKT 点可能是鞍点 / 局部极大 \\
“共轭两次必回原函数” & 仅闭凸函数 $f^{\ast\ast}=f$，否则得闭凸包 \\
“越新的算法越好” & 算法只有 “更适合”；先识别结构再查收敛率 \\
\bottomrule
\end{tabular}
\end{table}

\section*{故障排查手册}
\begin{enumerate}
  \item \textbf{症状}：CVXPY 报 “problem does not follow DCP rules”。\textbf{原因}：表达式凸但不满足 DCP 合成规则（编译器无法自动验证，如 \texttt{sqrt(quad\_form(x,P))}）。\textbf{对策}：用 DCP-compatible 的 atom（如 \texttt{norm(P\_sqrt@x)}）或手动等价变换（\cref{subsec:cvx-fun-ops}）。
  \item \textbf{症状}：求解器报 “primal infeasible” 但问题应当可行。\textbf{排查}：检查是否 Slater / 约束规范失败（\cref{subsec:cvx-slater}）——退化的等式 / 边界相切会使对偶变量不存在。
  \item \textbf{症状}：梯度下降收敛极慢或振荡发散。\textbf{原因}：步长 $1/L$ 中 $L$ 估错（$L=\lambda_{\max}(\nabla^2 f)$）或问题病态（$\kappa$ 大）。\textbf{对策}：用 Armijo 回溯自适应步长；或换加速 / 预处理（\cref{ins:condition-number}）。
  \item \textbf{症状}：ADMM 几千步不收敛。\textbf{原因}：罚参数 $\rho$ 未按问题尺度初始化。\textbf{对策}：用残差平衡自适应 $\rho$；OSQP 默认做 Ruiz 均衡。
  \item \textbf{症状}：手算 KKT 解出 $\lambda_i^\star<0$。\textbf{原因}：误设了 active 约束集，或符号约定不一致。\textbf{对策}：先解无约束 / 部分 active 情形检验，再逐步加 active 约束（\cref{thm:kkt} 证明后的数值流程）。
  \item \textbf{症状}：用 Fenchel--Moreau（$f^{\ast\ast}=f$）得到的不是原函数。\textbf{原因}：$f$ 非闭凸（非 proper 或非 lsc）。\textbf{对策}：先验证 proper $+$ lsc $+$ convex（\cref{thm:fenchel-moreau}），自定义分段函数尤需检查 $\operatorname{epi}f$ 是否闭。
\end{enumerate}
